\documentclass{article}

%% AMS packages
\usepackage{amsfonts,amsmath,amstext,amsbsy,amssymb}
\usepackage{amsopn,amsthm,amscd,amsxtra}

\usepackage{dsfont}

%% Graphics
\usepackage{graphicx}

%% Figures
% \usepackage{epsfig} \input{xy} \xyoption{all}

 \usepackage{url}


%% LATeX packages
\usepackage{latexsym,mathrsfs}



\usepackage{thmtools}
\usepackage{setspace}
\usepackage{geometry}
\usepackage{float}
\usepackage{hyperref}
\usepackage[utf8]{inputenc}
\usepackage[english]{babel}
\usepackage{framed}
\usepackage[dvipsnames]{xcolor}
\usepackage{tcolorbox}


%% Special arrows
\usepackage{mathabx} \newcommand{\carr}{\righttoleftarrow}
\newcommand{\wto}{\rightharpoonup}

%% Hyperlinks setup
\usepackage{hyperref}
\RequirePackage[dvipsnames]{xcolor} % [dvipsnames]
\definecolor{halfgray}{gray}{0.55}
% chapter numbers will be semi transparent .5 .55 .6 .0
\definecolor{webgreen}{rgb}{0,0.5,0}
\definecolor{webbrown}{rgb}{.6,0,0} \hypersetup{%
  colorlinks=true, linktocpage=true, pdfstartpage=3,
  pdfstartview=FitV,%
  breaklinks=true, pdfpagemode=UseNone, pageanchor=true,
  pdfpagemode=UseOutlines,%
  plainpages=false, bookmarksnumbered, bookmarksopen=true,
  bookmarksopenlevel=1,%
  hypertexnames=true,
  pdfhighlight=/O,%hyperfootnotes=true,%nesting=true,%frenchlinks,%
  urlcolor=webbrown, linkcolor=RoyalBlue,
  citecolor=webgreen, %pagecolor=RoyalBlue,%
  % uncomment the following line if you want to have black links
  % (e.g., for printing)
  % urlcolor=Black, linkcolor=Black,
  % citecolor=Black, %pagecolor=Black,%
  pdftitle={},%
  pdfauthor={},%
  pdfsubject={2000 Mathematical Subject Classification: Primary:},%
  pdfkeywords={},%
  pdfcreator={pdfLaTeX},%
  pdfproducer={LaTeX with hyperref}%
}

%%%%%%%%%%%%%%%%%%%%%%%%%%%%%%

%% Portuguese
% \usepackage[brazilian]{babel}
\usepackage{ae}

%%%%%%%%%%%%%%%%%%%%%%%%%%%%%%

% Bracket-like operations
\newcommand{\floor}[1]{\left\lfloor{#1}\right\rfloor}
\newcommand{\ceil}[1]{\left\lceil{#1}\right\rceil}
\newcommand{\abs}[1]{\left\lvert{#1}\right\rvert}
\newcommand{\norm}[1]{\left\|{#1}\right\|}
\newcommand{\gen}[1]{\left\langle{#1}\right\rangle}
\newcommand{\scprod}[2]{\left\langle{#1},{#2}\right\rangle}

% Typefaces
\newcommand{\bb}{\mathbb} \newcommand{\mbf}{\mathbf}
\newcommand{\mc}{\mathcal} \newcommand{\ms}{\mathscr}
\newcommand{\ds}{\mathds}

% Usual spaces
\newcommand{\R}{\mathbb{R}}\newcommand{\N}{\mathbb{N}}
\newcommand{\Z}{\mathbb{Z}}\newcommand{\Q}{\mathbb{Q}}
\newcommand{\T}{\mathbb{T}}\newcommand{\C}{\mathbb{C}}
\newcommand{\D}{\mathbb{D}}\newcommand{\A}{\mathbb{A}}

%% Abbreviations
\newcommand{\cf}{c.f.\ } \newcommand{\ie}{i.e.\ }
\newcommand{\eg}{e.g.\ } \newcommand{\etal}{\textit{et al\ }}


%% Diff-like operators
\newcommand{\Hom}{\mathrm{Hom}}
\newcommand{\Diff}[2]{\mathrm{Diff}_{#1}^{#2}}
\newcommand{\Homeo}[1]{\mathrm{Homeo}_{#1}}
\newcommand{\Dis}[1]{\mathcal{D}^\prime_{#1}}


%% Measures and differentials
\newcommand{\Leb}{\mathrm{Leb}} \newcommand{\dd}{\:\mathrm{d}}

%% Some operators
\DeclareMathOperator{\M}{\mathfrak{M}} \DeclareMathOperator{\Lip}{Lip}
\DeclareMathOperator{\Per}{Per} \DeclareMathOperator{\Fix}{Fix}
\DeclareMathOperator{\SL}{SL} \DeclareMathOperator{\GL}{GL}
\DeclareMathOperator{\spec}{spec}
\DeclareMathOperator{\rank}{rank} \DeclareMathOperator{\diver}{div}
\DeclareMathOperator{\cl}{cl} \DeclareMathOperator{\sign}{sign}
\DeclareMathOperator{\diam}{diam}

\newcommand{\cc}{\mathrm{cc}} % connected components

\newcommand{\pr}[1]{\mathrm{pr}_{#1}} % projections


%% For marginal notes
\newcommand{\marginal}[1]{\marginpar{{\scriptsize {#1}}}}

%% Topological dynamics on tori
\newcommand{\Symp}[1]{\mathrm{Symp}_{#1} (\mathbb{T}^2)}
\newcommand{\Ham}{\mathrm{Ham} (\mathbb{T}^2)}
\newcommand{\Thomeo}{\widetilde{\mathrm{Homeo}_0}}
\newcommand{\Tsymp}{\widetilde{\mathrm{Symp}_0} (\mathbb{T}^2)}
\newcommand{\Tham}{\widetilde{\mathrm{Ham}} (\mathbb{T}^2)}
\newcommand{\wh}{\widehat}


%% Especial notation for this article

\newcommand{\B}{\mathscr{B}}
\newcommand{\Bs}{\mathcal{S}}
\newcommand{\eps}{\varepsilon}

%%% Settings of original template

\colorlet{LightGray}{White!90!Periwinkle}
\colorlet{LightOrange}{Orange!15}
\colorlet{LightGreen}{Green!15}
\colorlet{LightYellow}{Yellow!15}
\colorlet{LightBlue}{Blue!15}

\newcommand{\HRule}[1]{\rule{\linewidth}{#1}}

\declaretheoremstyle[name=Theorem,]{thmsty}
\declaretheorem[style=thmsty,numberwithin=section]{theorem}
\tcolorboxenvironment{theorem}{colback=LightOrange}

\declaretheoremstyle[name=Corollary,]{corsty}
\declaretheorem[style=corsty,numberlike=theorem]{corollary}
\tcolorboxenvironment{corollary}{colback=LightOrange}


\declaretheoremstyle[name=Proposition,]{prosty}
\declaretheorem[style=prosty,numberlike=theorem]{proposition}
\tcolorboxenvironment{proposition}{colback=LightYellow}

\declaretheoremstyle[name=Lemma,]{lemsty}
\declaretheorem[style=lemsty,numberlike=theorem]{lemma}
\tcolorboxenvironment{lemma}{colback=LightGreen}

\declaretheoremstyle[name=Definition,]{defsty}
\declaretheorem[style=defsty,numberwithin=section]{definition}
\tcolorboxenvironment{definition}{colback=LightBlue}


\declaretheoremstyle[name=Problem,]{probsty}
\declaretheorem[style=probsty,numberwithin=section]{problem}
\tcolorboxenvironment{problem}{colback=LightGray}

\declaretheoremstyle[name=Example,]{examsty}
\declaretheorem[style=exasty,numberwithin=section]{example}
\tcolorboxenvironment{example}{colback=LightBlue}


\declaretheoremstyle[name=Exercise,]{exesty}
\declaretheorem[style=exesty,numberlike=problem]{exercise}
\tcolorboxenvironment{exercise}{colback=LightGray}



\setstretch{1.2}
\geometry{
    textheight=9in,
    textwidth=5.5in,
    top=1in,
    headheight=12pt,
    headsep=25pt,
    footskip=30pt
}


 %% Fancy indexing in enumerations
\usepackage{enumerate}

%% Figures
% \usepackage{epsfig} \input{xy} \xyoption{all}


\setstretch{1.2}
\geometry{
    textheight=9in,
    textwidth=5.5in,
    top=1in,
    headheight=12pt,
    headsep=25pt,
    footskip=30pt
}

% ------------------------------------------------------------------------------

\begin{document}

% ------------------------------------------------------------------------------
% Cover Page and ToC
% ------------------------------------------------------------------------------

\title{ \normalsize \textsc{Doutorado em Matemática - IMPA}
        \\ [2.0cm]
        \HRule{1.5pt} \\
        \LARGE \textbf{\uppercase{Lecture notes on Ergodic Theory}
        \HRule{2.0pt} \\ [0.6cm] \LARGE{} \vspace*{10\baselineskip}}}
\date{}
\author{Alejandro Kocsard \\
        IMPA\\
        \today}

\maketitle
\newpage

\tableofcontents
\newpage

% ------------------------------------------------------------------------------
% Chapters
% ------------------------------------------------------------------------------

\section{Dynamical Systems}
\label{sec:dynamical-systems}

In this course we shall deal with \emph{autonomous deterministic dynamical systems.} They can be used to represent or modelize any phenomenon that changes along the time in such a way that the current state of the system completely determines the future evolution and the law that drives this evolution does not change along the time. Mathematically, an \textbf{autonomous deterministic dynamical system} (a \textbf{dynamical system}, for short) is given by a set $M$, usually called the \textbf{phase space,} which represents the set of possible states of the system, a subsemigroup $\T\subset\R$ that represents the \textbf{time set,} and \emph{``a law''} $\Phi\colon\T\times M\to M$ that determines the evolution of the system in such a way that given the current state of our system $x_{0}\in M$, at time $t\in\T$ the corresponding state will be $x_{t}=\Phi(x_{0},t)$. The autonomy of the system is reflected in the so called \textbf{flow} (or \textbf{semi-flow}) condition function $\Phi$ satisfies:
\begin{equation}\label{eq:flow-condition}
  \Phi(t+s,x)=\Phi\big(t,\Phi(s,x)\big) = \Phi\big(t,\Phi(s,x)\big), \quad\forall t,s\in\T,\ \forall x\in M.
\end{equation}

We will consider two different kinds of dynamical systems: \textbf{discrete time} and \textbf{continuous time} ones, \ie~where $\T=\N_{0}$ and $\T=\R_{\geq 0}$, respectively. Moreover, we will also consider \textbf{invertible systems} where the current state of the systems does not just determine the future but the past as well. When dealing with invertible discrete dynamical systems, the time set will be $\T=\Z$, and in the case of invertible continuous time ones, we will have $\T=\R$.

The main goal of Dynamical Systems consists in describing the \textbf{asymptotic behavior of the orbits of the systems,} \ie~the sets
\begin{displaymath}
  \mc{O}_{\Phi}(x):=\left\{\Phi(t,x) : t\in\T\right\}, \quad\forall x\in M.
\end{displaymath}


\begin{exercise}[Discrete time generator]
  Given a discrete time dynamical system $\Phi\colon\T\times M\to M$, show that $\Phi$ satisfies the flow condition~\eqref{eq:flow-condition} if and only if there is $f\colon M\carr$ such that
  \begin{displaymath}
    \Phi(n,x)=f^{n}(x), \quad\forall x\in M,\ \forall n\in\N_{0},
  \end{displaymath}
  where the iteration of $f$ is inductively defined by $f^{0}:=id_{M}$ and $f^{n+1}=f\circ f^{n}$, for every $n\geq 0$. Show that $\Phi$ is invertible if and only if $f$ is bijective and, in such a case, one can define $f^{-n}:={\big(f^{-1}\big)}^{n}$, for all $n\geq 0$.
\end{exercise}

On the other hand, one can find \emph{infinitesimal generators} for continuous time smooth dynamical systems:

\begin{exercise}[Continuous time generator]
  Let $M$ be a differentiable closed manifold (\ie~compact and boundaryless) and $\Phi\colon\R\times M\to M$ a continuous time dynamical system (that satisfies the flow condition~\ref{eq:flow-condition}). If  we assume $\Phi$ is a $C^{2}$ map and we define $X(p):=\partial_{t}\Phi(t,p)\big|_{t=0}\in T_{p}M$, for every $p\in M$, then the function $\Phi(\cdot,x_{0})\colon\R\to M$ is the unique solution of the \emph{initial value problem}
  \begin{displaymath}
    \dot x(t)=X\big(x(t)\big), \quad x(0)=x_{0},
  \end{displaymath}
  for every $x_{0}\in M$. In other words, the vector field $X$ is the infinitesimal generator of the flow $\Phi$.
\end{exercise}

To give meaning to the expression \emph{``asymptotic behavior''} it is necessary to endow dynamical systems with an additional structure that makes it possible to analyze the behavior of their orbits as time tends to infinity.

When the phase space $M$ is a topological space and the semi-flow $\Phi$ is continuous, one can study the accumulation points of the orbits as time goes to infinity and this is the basic scenario of \textbf{topological dynamics.}

On the other hand, when $\Phi\colon\T\times M\to M$ is an arbitrary dynamical system, one can study the \textbf{statistical properties} the frequency of visits of the orbit of a certain point $x\in M$ to a subset $A\subset M$. This problem naturally lead us to consider \textbf{invariant measures} on $M$ and this is one of the fundamental problems in \textbf{ergodic theory.}

\subsection{The setup of Ergodic Theory}\label{sec:setup-ergodic-theory}

Ergodic theory can be roughly described as the study of \textbf{measure preserving dynamical systems.}

More precisely, by a \textbf{(discrete time) dynamical system} we mean a map $T\colon X\to X$, where $X$ is its \textbf{phase space} and the dynamics is given by the iteration of $T$: $T^{0}=id_{X}$ and $T^{n+1}:= T\circ T^{n}$, for every $n\geq 0$. When $T$ is bijective, one can consider the ``past trajectories'' taking negative iterates that are defined by $T^{-n}:={(T^{-1})}^{n}$, for every $n\geq 0$.

We say that $T\colon X\carr$ is a \textbf{measurable} system when $X$ is a endowed with a $\sigma$-algebra $\B\subset 2^{X}$ and $T$ is a $\B$-measurable map, \ie~$T^{-1}(B)\in\B$, for every $B\in\B$.

When $\mu$ is a measure on the measurable space $(X,\B)$ and $T\colon X\carr$ is a measurable map, we say that $T$ \textbf{preserves} $\mu$, or that $\mu$ is a \textbf{$T$-invariant measure}, when it holds
\begin{displaymath}
  \mu(A)=\mu\big(T^{-1}(A)\big),\quad\forall A\in\B.
\end{displaymath}

So, a \textbf{measure preserving system} is the quadruple $T\colon (X,\B,\mu)\carr$ where $T\colon X\carr$ is a $\B$-measurable map and $\mu$ is a $T$-invariant \textbf{probability} measure, \ie~where $\mu(X)=1$.

Analogous concepts can be defined for continuous time dynamical systems, also called \textbf{semi-flows} and \textbf{flows}, and more generally, for (measurable) actions of semi-groups and groups.

Our main references for the course are the books of Viana and Oliveira~\cite{VianaOliveiraFoundErgTh} and Mañé~\cite{ManeBook}.

\subsection{Motivations}\label{sec:motivations} Let us see some fundamental problems in Dynamical Systems that motivate the study of invariant measures and their properties.

\subsubsection{Dynamical systems with a natural invariant measure}\label{sec:dynam-syst-natur-inv-meas}

There are many classical dynamical systems, which are used to model natural phenomena, that already arise with a ``natural'' invariant (finite) measure. From the historical point of view, one of the most important examples are Hamiltonian systems that are used to described the evolution of conservative systems in Classical Mechanics.

So, the first main question in ergodic theory could be briefly and roughly summarized as follows:
\begin{quote}
  \emph{What dynamical information can be extracted from the fact that a dynamical system leaves invariant a certain measure?}
\end{quote}

\subsubsection{Statistical study of dynamical systems}\label{sec:statical-study-dyn-sys}

A second motivation to study measure preserving systems comes from the analysis of \emph{statistical aspects} of dynamical systems. In its simplest form, let us suppose a dynamical system $T\colon X\carr$ is given and there is a point $x\in X$ such that \emph{asymptotic visiting frequency function of $x$} is well defined for every $A\in 2^x:=\{B\subset X\}$. More precisely, let us suppose that defining $\tau_{x}^{(n)}\colon 2^X\to [0,1]$ by
\begin{displaymath}
  \tau_{x}^{(n)}(A):=\frac{\sharp\left\{0\leq i\leq n-1 : T^{i}(x)\in A\right\}}{n}, \quad\forall A\subset X,\ \forall n\in\N,
\end{displaymath}
the following limit exists:
\begin{displaymath}
  \tau_x(A):=\lim_{n\to+\infty}\tau_x^{(n)}(A), \quad\forall A\subset X.
\end{displaymath}

\begin{exercise}
  Show that $\tau_x\colon 2^X\to[0,1]$ is a probability finitely additive measure and it holds
  \begin{displaymath}
    \tau_x\big(T^{-1}(A)\big)=\tau_x(A), \quad\forall A\subset X.
  \end{displaymath}

  Is $\tau_x$ $\sigma$-additive?
\end{exercise}

As we shall see along the course, in general it is not true that the asymptotic visiting frequency is well defined for every set $A$ and different points $x$c an have different asymptotic visiting frequencies on the same set $A$. So, in order to study the statistical properties of a certain system it will be necessary to analyze the whole set of invariant measures.


\subsubsection{Applications to combinatorics and number theory}\label{sec:appl-combinatorics-number-theory}

As third motivation, we will see that ideas and techniques of ergodic theory can be applied to get some nice results on combinatorics and number theory. The main reference for this part of the course in the book of Einsiedler and Ward~\cite{EinsiedlerWardErgThViewNumbTh}.


\subsubsection{Normal numbers}\label{sec:normal-numbers}

A real number $x\in [0,1]$ is said to be \emph{$b$-normal}, for $b\in\N$ and $b\geq 2$, when its digit in base $b$ are uniformly distributed, \ie~given a sequence of integer numbers $(a_n)_{n\geq 1}$, with $0\leq a_n\leq b-1$ such that
\begin{displaymath}
  x=\sum_{n\geq 1} \frac{a_n}{b^n},
\end{displaymath}
it holds
\begin{displaymath}
  \lim_{n\to+\infty}\frac{\sharp\{i\in\N : i\leq n-1,\ a_i=a\}}{n}=\frac{1}{b},
\end{displaymath}
for every $a\in\{0,1,\ldots,b-1\}$.

Then we say that $x\in[0,1]$ is just \emph{normal} when it is $b$-normal, for every $b\in\N$ with $b\geq 2$.

Using ergodic techniques one can prove that Lebesgue almost every point of $[0,1]$ is normal indeed.


\subsubsection{Distribution of the first digit of $2^n$}\label{sec:distr-firs-dig-2n}

For each $n\in\N$, let us writ $a_n$ for the first digit of the decimal representation of $2^n$. We want to study the asymptotic distribution of the sequence $(a_n)$ in the set $\{1,\ldots,9\}$.

In this case, using ergodic tools again, one can show that
\begin{displaymath}
  \frac{\sharp\left\{1\leq j\leq n : a_j=k\right\}}{n} \to\log_{10}\left(\frac{k+1}{k}\right), \quad\text{as } n\to+\infty,
\end{displaymath}
for each $k\in\{1,\ldots,9\}$.

\subsubsection{Existence of arithmetic progressions}\label{sec:exist-arit-progr}

One says the set $A\subset\Z$ has \emph{positive upper density} when there are sequences of integer numbers $(m_j)$ and $(n_j)$ such that $m_j-n_j\to+\infty$, as $j\to+\infty$ and
\begin{displaymath}
  \lim_{j\to+\infty} \frac{\sharp(A\cap [m_j,n_j])}{m_j-n_j} >0.
\end{displaymath}

Erd\H{o}s and Turán proposed the following conjecture in~\cite{ErdosTuranSeqIntegers}: Every subset of integer numbers with positive upper density contains arbitrarily long arithmetic progressions. This conjecture remained open for almost 40 years and it was finally proved by Szeméredi~\cite{SzemerediSetsIntContNoKElemArithProg} using rather complicated combinatorial arguments. Two years later, Furstenberg proposed in~\cite{FurstenbergErgBehavDiagMeasSzemeredi} a purely ergodic proof of this result, opening the door to a long standing and fruitful interaction of Combinatorics and Ergodic Theory.


\section{Measure theory and the setting of Ergodic theory}\label{sec:measure-theory}

\subsection{Measurable spaces}\label{sec:measurable-spapces}

All along these notes, $X$ will denote a non-empty set. We shall write $2^{X}$ for the set of subsets of $X$. A \textbf{$\sigma$-algebra} on $X$ is a family $\mc{B}\subset 2^{X}$ such hat $\emptyset,X\in\mc{B}$ and given $A_{1},A_{2},\ldots\in\mc{B}$, it holds $\bigcup_{n\geq 1} A_{n}\in\mc{B}$. A \textbf{measurable space} is a pair $(X,\mc{B})$ where $X$ is a non-empty set and $\mc{B}$ is a $\sigma$-algebra on $X$.

There is a natural partial order among $\sigma$-algebras of $X$ given by set inclusion. Invoking this order, given any family $\mc{C}\subset 2^{X}$ one can define the \textbf{$\sigma$-algebra generated by $\mc{C}$}, which will be denoted by $\sigma({\mc{C}})$, as the smallest $\sigma$-algebra that contains $\mc{C}$.

When $X$ is a topological space and $\tau\subset 2^{X}$ is the topology of $X$, \ie~$\tau$ is the family of open subsets in $X$, the $\sigma$-algebra $\sigma(\tau)$ generated by $\tau$ is called the \textbf{Borel $\sigma$-algebra} of $X$ is usually denoted by $\mc{B}(X)$.


Given a measurable space $(X,\mc{B})$, and an arbitrary map $f\colon X\to Y$, one can define the \textbf{push-forward $\sigma$-algebra} $f_{\star}\mc{B}$ on $Y$ by
\begin{equation}
  \label{eq:push-forward-sigma-alg-def}
  f_{\star}\mc{B}:=\left\{E\subset Y : f^{-1}(E)\in\mc{B}\right\}.
\end{equation}

On the other hand, given another map $h\colon Z\to X$, we defined the \textbf{pull-back $\sigma$-algebra} $f^{\star}\mc{B}$ on $Z$ by
\begin{equation}
  \label{eq:pull-back-sigma-alg-def}
  f^{\star}\mc{B}:=\left\{E\subset Z : \exists A\in\B,\ f^{-1}(A)=E\right\}.
\end{equation}

Given two measurable spaces $(X,\mc{B})$ and $(Y,\mc{C})$, a map $f\colon X\to Y$ is said to be \textbf{measurable} when $\mc{C}\subset f_{\star}\mc{B}$.

\begin{exercise}
  If $(X,\mc{B})$ and $(Y,\mc{C})$ are measurable spaces and $f\colon X\to Y$, show that $f^{\star}\mc{C}\subset\mc{B}$ if and only if $\mc{C}\subset f_{\star}\mc{B}$.
\end{exercise}

A \textbf{measure} on a measurable space $(X,\mc{B})$ is map $\mu\colon\mc{B}\to[0,+\infty]$ such that $\mu(\emptyset)=0$, $\mu(X)>0$ and
\begin{displaymath}
  \mu\bigg(\bigsqcup_{n=1}^{\infty} A_{n}\bigg) = \sum_{n=1}^{\infty} \mu(A_{n}),
\end{displaymath}
for every sequence of pairwise disjoint sets $A_{1},A_{2},\ldots \in\mc{B}$.
The measure $\mu$ on $(X,\mc{B})$ is said to be \textbf{finite} when $\mu(X)<\infty$; a \textbf{probability} measure when $\mu(X)=1$; and \textbf{$\sigma$-finite} when there exists a sequence $X_{1},X_{2},\ldots\in\mc{B}$ such that $\mu(X_{n})<\infty$ for every $n\geq 1$ and $X=\bigcup_{n\geq 1} X_{n}$.

Given a measure space $(X,\mc{B},\mu)$, we say that two sets $A,B\in\mc{B}$ are \textbf{equal modulo $\mu$} when $\mu(A\triangle B)=0$, where $A\triangle B:=(A\setminus B)\cup (B\setminus A)$.

\begin{exercise}
  Let $(X,\mc{B},\mu)$ be a finite measure space. Show that the relation $\sim$ on $\mc{B}$ given by $A\sim B$ if and only if $A$ and $B$ are equal modulo $\mu$ is an equivalence relation. Then, let us consider the function $d_{\mu}\colon \mc{B}/\!\!\sim\times \mc{B}/\!\!\sim\to\R$ given by
  \begin{displaymath}
    d_{\mu}([A],[B]):=\mu(A\triangle B), \quad\forall A,B\in\B,
  \end{displaymath}
  and where $[A],[B]$ denotes the equivalence classes of $A$ and $B$, respectively. Show that the function $d_{\mu}$ is well defined and is a distance on $\mc{B}/\!\!\sim$.
\end{exercise}



A \textbf{measure space} is triple $(X,\mc{B},\mu)$, where $(X,\mc{B})$ is a measurable space and $\mu$ is a measure on it. The $\sigma$-algebra $\mc{B}_{0}^{\mu}\subset\mc{B}$ of \textbf{$\mu$-null sets} is given by
\begin{displaymath}
  \mc{B}_{0}^{\mu}:=\{N\in\mc{B} : \mu(N)=0\}.
\end{displaymath}

We say that the measure space $(X,\mc{B},\mu)$ is \textbf{complete} when for every $N\in\mc{B}_{0}^{\mu}$ and any $N'\subset N$, it holds $N\in\mc{B}$.

Every measure space $(X,\mc{B},\mu)$ admits a unique \textbf{completion:} it is a complete metric space $(X,\overline{\mc{B}},\bar\mu)$ such that
\begin{displaymath}
  \overline{\mc{B}}:=\left\{B\cup N_{1}\setminus N_{2} : B\in\mc{B},\ \exists N\in\mc{B}_{0}^{\mu},\ N_{1},N_{2}\subset N\right\},
\end{displaymath}
and
\begin{displaymath}
  \bar{\mu}(B\cup N_{1}\setminus N_{2}):=\mu(B), \quad\forall B\in\mc{B},\ \forall N\in\mc{B}_{0}^{\mu},\ \forall N_{1},N_{2}\subset N.
\end{displaymath}

So, in order to avoid unnecessary technical difficulties, from now on we will assume all our measure spaces are complete.

Given a measure space $(X,\mc{B},\mu)$ and a map $f\colon X\to Y$, we can define the \textbf{push-forward measure} $f_{\star}\mu$ on the measurable space $(Y,f_{\star}\mc{B})$ by
\begin{equation}
  \label{eq:push-forward-measure-def}
  f_{\star}\mu(A):=\mu\big(f^{-1}(A)\big), \quad\forall A\in f_{\star}\mc{B}.
\end{equation}

\begin{exercise}\label{exer:pull-back-function-push-forward-measures}
  Let $(X,\mc{B},\mu)$ be a measure space, $(Y,\mc{C})$ a measurable space and $f\colon X\to Y$ a measurable map. The it holds
  \begin{displaymath}
    \int_{X}\phi\circ f\dd\mu = \int_{Y}\phi\dd f_{\star}\mu,
  \end{displaymath}
  for any measurable function $\phi\colon Y\to [0,\infty]$.
\end{exercise}

It is important to remark that, in general, it is not possible to pulled measures back. In fact, it is only possible to do it when the corresponding map is invertible.

Given two measure spaces $(X,\mc{B},\mu)$ and $(Y,\mc{C},\nu)$, a \textbf{measure preserving map} or a \textbf{metric morphism} is a measurable map $h\colon X\to Y$ such that $h_{\star}\mu=\nu$. Observe that if we were pedantically formal, taking into account that the domain of $h_{\star}\mu$ is $h_{\star}\mc{B}\supset\mc{C}$, last equality should be written as $h_{\star}\mu\big|_{\mc{C}}=\nu$. In particular, a map $f\colon X\to X$ is said to be a \textbf{metric endomorphism} or a measure preserving map when $f$ is measurable and $f_{\star}\mu=\mu$. In such a case, we say that $\mu$ is an \textbf{$f$-invariant measure.}

A measure preserving map $h\colon (X,\mc{B},\mu)\to (Y,\mc{C},\nu)$ is said to be a \textbf{measure isomorphism} when there exist measurable sets $X'\in\mc{B}$ and $Y'\in\mc{C}$ such that $\mu(X\setminus X')=\nu(Y\setminus Y')=0$, $h\big|_{X'}\colon X'\to Y'$ is a bijection and $h^{-1}\big|_{Y'}$ is a measurable map.

Finally, we will say that a measure preserving map $f\colon (X,\mc{B},\mu)\to (X,\mc{B},\mu)$ is an \textbf{automorphism} or an \textbf{invertible measure preserving map} when $f$ is a metric isomorphism.

\subsection{Borel and Lebesgue spaces}\label{sec:Borel-Lebesgue-spaces}

Let $M$ denote a separable metric space and let $\B_{M}$ be its Borel $\sigma$-algebra, \ie~$\B_{M}$ is the smallest $\sigma$-algebra containing every open subset.

A probability space $(X,\B,\mu)$ is said to be a \textbf{Borel space} when there is a Borel probability measure $\mu_{M}$ on $(M,\B_{M})$ such that $(X,\B,\mu)$ is isomorphic to $(M,\B_{M},\mu_{M})$.

Observe that if $(X,\B,\mu)$ is Borel space, then there exists a countable family $\mc{A}\subset\B$ such that $\sigma(\mc{A})=\B$, \ie~the $\sigma$-algebra $\B$ is countably generated.

Finally, we say that the probability space $(X,\B,\mu)$ is a \textbf{Lebesgue space} when it is complete and is isomorphic to the completion of a Borel probability space. See~\cite[Appendix A. 6]{EinsiedlerWardErgThViewNumbTh} for more details about Borel and Lebesgue spaces.



\subsection{Smooth measures}\label{sec:smooth-measures}

Let $M$ and $N$ be a smooth manifolds, $\alpha$ be a $k$-form on $N$ and $f\colon M\to N$ be a $C^{1}$ map. The \textbf{pullback} of $\alpha$ by $f$ is the $k$-form $f^{\star}\alpha$ on $M$ given by
\begin{displaymath}
  f^{\star}\alpha(X_{1},X_{2},\ldots,X_{k}):= \alpha\big(Df_{x}(X_{1}),Df_{x}(X_{2}),\ldots,Df_{x}(X_{k})\big),
\end{displaymath}
for all $x\in M$, and every $X_{1},\ldots,X_{k}\in T_{x} M$.

On the other hand, given a complete smooth vector field $X$ on the manifold $M$, we can consider the flow $\Phi_{X}\colon\R\times M\to M$ induced by $X$ on $M$. We say that a Borel measure $\mu$ on $M$ is $X$-invariant (or $\Phi_{X}$-invariant) when it is $\Phi_{X}^{t}$-invariant, for every $t\in\R$, where $\Phi_{X}^{t}(z)=\Phi_{X}(z,t)$, for every $z\in M$ and any $t\in\R$.

Then, given a smooth $k$-form $\alpha$ on $M$, one defines the \emph{Lie derivative} of $\alpha$ with respect to $X$ by
\begin{equation}
  \label{eq:Lie-derivative-def}
  \mathcal{L}_{X}\alpha := \lim_{t\to 0} \frac{{\big(\Phi_{X}^{t}\big)}^{\star}\alpha-\alpha}{t}.
\end{equation}

Now, given $M$ a compact oriented smooth $d$-manifold $M$, a $d$-form $\omega$ is said to be \textbf{volume form}, when for every $x\in M$ and every positively oriented basis $\{X_{1}, X_{2},\ldots , X_{d}\}$ of the tangent space $T_{x}M$, it holds $\omega_x(X_{1},\ldots,X_{d})>0$.

Then, we can define the continuous linear functional $\Lambda_{\omega}\colon C^{\infty}(M,\R)\to\R$ by
\begin{displaymath}
  \Lambda_{\omega}(\phi):=\int \phi\omega, \quad\forall \phi\in C^{\infty}(M,\R).
\end{displaymath}
One can easily check that $\Lambda_{\omega}$ can be uniquely extended to a continuous linear functional $\overline{\Lambda_{\omega}}\colon C^{0}(M,\R)\to\R$, and by Riesz' representation theorem we conclude that there exists a finite Borel measure $\mu_{\omega}$ such that
\begin{displaymath}
  \overline{\Lambda_{\omega}}(\phi) = \int_{M}\phi\dd\mu_{\omega}, \quad\forall \phi\in C^{0}(M,\R).
\end{displaymath}

The measure $\mu_{\omega}$ is said to be a \textbf{smooth measure.}
\begin{exercise}
  Given a $C^1$-map $f\colon M\carr$, show that the smooth measure $\mu_\omega$ is $f$-invariant if and only if $f^\star\omega=\omega$.

  On the other hand, given a vector field $X$ on $M$, show that $\mu_\phi$ os $X$-invariant if and only if $\mc{L}_X\omega\equiv 0$.
\end{exercise}

\begin{exercise}
  Can we define smooth measures on non-compact oriented manifolds? Can smooth measures be finite on non-compact manifolds? Smooth measures can be defined on non-orientable manifolds?
\end{exercise}

\subsection{Some important exercises from~\cite[\S1.1.1]{VianaOliveiraFoundErgTh}}\label{sec:some-import-exerc-1.1.1}

1.1.2, 1.1.3, 1.1.4, 1.1.5.


\section{Recurrence on finite measure spaces}\label{sec:recurrence-Poincare}

In this section we shall analyze the fact that finiteness of invariant measures imposes recurrence of the dynamical system.

Our first result is the celebrated

\begin{theorem}[Poincaré recurrence theorem]\label{thm:poincare-recurrence-v-return-sets-for-1-time}
  Let $(X,\mc{B},\mu)$ be a probability space and $f\colon (X,\mc{B},\mu)\carr$ be a measure preserving map. Let $A\in\B$ be a measurable set such that $\mu(A)>0$. Then, there exists $n\geq 1$ such that
  \begin{displaymath}
    \mu\big(A\cap f^{-n}(A)\big)>0.
  \end{displaymath}
\end{theorem}

\begin{proof}
  Let us suppose that $A,f^{-1}(A),f^{-2}(A),\ldots$ are essentially disjoint, \ie~suppose that
  \begin{equation}
    \label{eq:fiA-fjA-disjoint-in-Poincare}
    \mu\big(f^{-i}(A)\cap f^{-j}(A)\big)=0, \quad\forall i,j\in\N_{0},\ i\neq j.
  \end{equation}

  Let us define the sets $A_{0}:=A$, and $A_{n}:=f^{-n}(A)\setminus \Big(\bigcup_{i=0}^{n-1} f^{-i}(A)\Big)$, for every $n\geq 1$. Then, the sets $A_{0},A_{1},A_{2},\ldots$ are pairwise disjoint, and by~\eqref{eq:fiA-fjA-disjoint-in-Poincare}, we have
  \begin{displaymath}
    \mu\big(f^{-n}(A)\big)\geq \mu(A_{n})\geq \mu(f^{-n}(A)) - \sum_{i=0}^{n-1} \mu\big(f^{-n}(A)\cap f^{-i}(A)\big) = \mu\big(f^{-n}(A)\big),
  \end{displaymath}
  for every $n\geq 1$, and hence it holds $\mu\big(f^{-n}(A)\big)=\mu(A_{n})$, for any $n\geq 0$. Since $\mu$ is $f$-invariant, this implies $\mu(A_{n})=\mu(A)>0$, for every $n\geq 0$, getting a contradiction because
  \begin{displaymath}
    1=\mu(X)\geq\mu\bigg(\bigsqcup_{i=0}^{\infty} A_{n}\bigg) = \sum_{i=0}^{\infty} \mu(A_{n}) = \sum_{i=0}^{\infty} \mu(A) = \infty.
  \end{displaymath}

  So, assumption~\eqref{eq:fiA-fjA-disjoint-in-Poincare} does not hold and hence, there exist $i,j\in\N$, with $i<j$ such that
  \begin{displaymath}
    0<\mu\big(f^{-i}(A)\cap f^{-j}(A)\big) = \mu\Big(f^{-i}\big(A\cap f^{i-j}(A)\big)\Big) = \mu(A\cap f^{-(j-i)}(A)).
  \end{displaymath}
\end{proof}

There are a lot of important results that can be proven as a consequence of Theorem~\ref{thm:poincare-recurrence-v-return-sets-for-1-time}. The firs one is indeed a very simple improvement of Theorem~\ref{thm:poincare-recurrence-v-return-sets-for-1-time}:

\begin{corollary}\label{cor:return-bounded-time-Poincare-thm}
  Let $f\colon(X,\mc{B},\mu)\carr$ and $A\in\mc{B}$ be as in Theorem~\ref{thm:poincare-recurrence-v-return-sets-for-1-time}. Then, there is a natural number $n\leq \frac{1}{\mu(A)}$ such that
  \begin{displaymath}
    \mu\big(A\cap f^{-n}(A)\big)>0.
  \end{displaymath}
\end{corollary}

\begin{exercise}
  Prove Corollary~\ref{cor:return-bounded-time-Poincare-thm}.
\end{exercise}


\begin{corollary}\label{cor:return-one-time-Poincare}
  Let $f\colon (X,\mc{B},\mu)\carr$ and $A\in\mc{B}$ as in Theorem~\ref{thm:poincare-recurrence-v-return-sets-for-1-time}. Then the set
  \begin{displaymath}
    A_{f}:=\left\{x\in A : \exists n\geq 1,\ f^{n}(x)\in A\right\}
  \end{displaymath}
  belongs to $\mc{B}$ and $\mu(A\setminus A_{f})=0$.
\end{corollary}

\begin{proof}
  First notice that
  \begin{displaymath}
    A_{f}=\bigcup_{n\geq 1} \big(A\cap f^{-n}(A)\big).
  \end{displaymath}
  So, $A_{f}\in\mc{B}$.

  Then, let us define $A':=A\setminus A_{f}$. Observe that $A'\cap f^{-n}(A')=\emptyset$, for every $n\geq 1$. Hence, we have that $f^{-m}(A')\cap f^{-n}(A')=\emptyset$, whenever $0\leq m<n$. By Theorem~\ref{thm:poincare-recurrence-v-return-sets-for-1-time}, we conclude that $\mu(A')=0$.
\end{proof}

Then notice that Corollary~\ref{cor:return-one-time-Poincare} can be considerably improved:

\begin{corollary}\label{cor:return-infinitely-many-times-Poincare}
  Let $f\colon (X,\mc{B},\mu)\carr$ and $A\in\mc{B}$ as in Theorem~\ref{thm:poincare-recurrence-v-return-sets-for-1-time}. Then the set
  \begin{displaymath}
    A_{\infty}:=\left\{x\in A : \exists n_{j}\uparrow\infty,\ f^{n_{j}}(x)\in A, \forall j\geq 1\right\}
  \end{displaymath}
  belongs to $\mc{B}$ and $\mu(A\setminus A_{\infty})=0$.
\end{corollary}

\begin{exercise}
  Prove Corollary~\ref{cor:return-infinitely-many-times-Poincare}.
\end{exercise}

\subsection{Topological recurrence}\label{sec:topological-recurrence-Poincare}

Let $X$ be a topological space and $f\colon X\carr$ be an arbitrary map. We say that a point $x_{0}\in X$ is \textbf{recurrent} for $f$ when given any neighborhood $V$ of $x_{0}$, there is a natural number $n\geq 1$ such that $f^{n}(x_{0})\in V$.

Recall that a \textbf{base} for the topology of $X$ is a family of open subsets $\mathcal{V}$ such that every open subset of $X$ can be writte as the union of elements of $\mc{V}$.

Then we have the following topological version of Poincaré recurrence theorem:

\begin{theorem}\label{thm:Poincare-recurrence-thm-topological-version}
  Let $X$ be a topological space admitting a countable base for its topology and let $\mc{B}(X)$ denote its Borel $\sigma$-algebra. Let $\mu$ be a Borel probability measure on $X$ (\ie~a probability measure on $\big(X,\mc{B}(X)\big)$ and $f\colon (X,\mc{B}(X),\mu)\carr$ a measure preserving transformation. Then, $\mu$-almost every point is recurrent for $f$.
\end{theorem}


\begin{proof}
  Let $\mathcal{U}:=\{U_{n} : n\in\N\}$ be a base for the topology of $X$. For each $n\in\N$ let us consider the set $W_{n}:=\{x\in U_{n}  : f^{k}(x)\not\in U_{n},\ \forall k\geq 1\}$. By Corollary~\ref{cor:return-one-time-Poincare} we know that $\mu(W_{n})=0$, for every $n\geq 1$. Hence, if we write $W:=\bigcup_{n\geq 1} W_{n}$, we have $\mu(W)=0$. Then we claim that every point in $X\setminus W$ is recurrent for $f$. In fact, let $x$ an arbitrary point of $X\setminus W$ and $V$ any neighborhood of $x$. Since $\mathcal{U}$ is a base for the topology of $X$, there is a natural number $n$ such that $x\in U_{n}\subset V$. Since $x\not\in W$, we have $x\not\in W_{n}$ and then, there is $k\geq 1$ such that $f^{k}(x)\in U_{n}\subset V$. Since $V$ was an arbitrary neighborhood, we conclude $x$ is recurrent for $f$, as desired.
\end{proof}

It is worth to recall at this point the following result due to Brikhoff:

\begin{theorem}[Birkhoff recurrence theorem]
  Let $(X,d)$ be a compact metric space and $f\colon X\carr$ be a continuous map. Then, there exists a recurrent point for $f$.
\end{theorem}

\begin{proof}
  Let us write $\mc{K}_f:=\{K\subset X : K \text{ is compact,}\ K\neq\emptyset,\ f(K)\subset K\}$. Then, $\mc{K}_f$ is clearly non-empty because $X\in\mc{K}_f$. On the other hand, if we consider the partial order relation $\subset$ given by inclusion, any totally ordered family of elements in $\mc{K}_f$ admits a lower bound in $\mc{K}_f$ which is given by the intersection of all the elements of the totally ordered family.

  So, as a consequence of Zorn's lemma, there is a minimal element $M$ in $\mc{K}_f$ with respect to the partial order $\subset$. One can easily check that given any $x\in M$, by minimality, its positive semi-orbit $\mc{O}_f^+(x):=\{f^n(x) : n\geq 0\}$ is dense in $M$, and consequently, $x$ is recurrent.
\end{proof}

\subsection{Some important exercises from~\cite[\S1.2.4]{VianaOliveiraFoundErgTh}}
\label{sec:some-import-exerc-1.2.4}

1.2.2, 1.2.3, 1.2.5, 1.2.6.


\section{Some examples of measure preserving transformations}\label{sec:examples-measure-pres-trans}

\subsection{Periodic orbits}\label{sec:examples-inv-meas-periodic-orbits}

Let $X$ be an arbitrary set and $z\in X$ an arbitrary point. Then we can define the $\delta_{z}$ on the measurable space $(X,2^{X})$ by
\begin{displaymath}
  \delta_{z}(A):=
  \begin{cases}
    1, &\text{if } z\in A;\\
    0, &\text{otherwise.}
  \end{cases}
\end{displaymath}


Now, if $f\colon X\carr$ is any map and there exists $p\in X$ which is \textbf{periodic} for $f$, \ie~there exists $k\in\N$ such that $f^{k}(z)=z$, then the \textbf{periodic measure}
\begin{displaymath}
  \nu_{p}:=\frac{1}{k}\sum_{j=0}^{k-1}\delta_{f^{j}(p)} = \frac{1}{k}\sum_{j=0}^{k-1} f_{\star}^{j}\delta_{p}
\end{displaymath}
is $f$-invariant.

\subsubsection{Homeomorphisms of the interval}\label{sec:homeos-interval}

Using periodic measures and Poincare's recurrence theorem (Theorem~\ref{thm:Poincare-recurrence-thm-topological-version}) one can easily determine all invariant measures for homeomorphisms of a compact interval\ldots at least for homeomorphisms with finitely many periodic orbits.

In fact, if $I\subset\R$ denotes a compact interval and $f\colon I\carr$ is an orientation-preserving homeomorphism (\ie~$f$ is an strictly increasing map), then one can easily show that $x\in I$ is a recurrent point for $f$ if and only if $f(x)=x$, and moreover, the set of fixed points $\Fix(f):=\{w\in I : f(w)=w\}$ is non-empty.

On the other hand, if $f$ is orientation-reversing, \ie~$f$ is strictly decreasing, the $\Fix(f)$ has exactly one point and any other recurrent point is periodic of period $2$.

In any case, if $f$ has finitely many periodic orbits, by Theorem~\ref{thm:Poincare-recurrence-thm-topological-version} one have that any $f$-invariant Borel probability measure is indeed a convex combination of periodic ones. What can we say about the characterization of $f$-invariant measures when $f$ has infinitely many periodic orbits?

\subsection{Expanding maps of the interval}\label{sec:expand-maps-interval}

Let $k\in\N$ be an arbitrary number and consider the map $f\colon [0,1]\carr$ given by
\begin{displaymath}
  f(x):=kx - \floor{kx}, \quad\forall x\in [0,1]
\end{displaymath}
where $\floor{x}:=\max\{\ell\in\Z : \ell\leq x\}$ denotes the \emph{integer part} of $x$.

We claim that the Lebesgue measure $m$ on $[0,1]$ is $f$-invariant. If $J\subset [0,1]$ is a subinterval, then one can easily check that $f^{-1}_{k}(J)$ is the disjoint union of $k$ intervals $J_{1},\ldots, J_{k}\subset [0,1]$ with $m(J_{i})=\frac{1}{k}m(J)$, for each $i\in\{1,\ldots, k\}$. So, we conclude
\begin{displaymath}
  m\big(f^{-1}(J)\big) = m\bigg(\bigsqcup_{i=1}^{k}J_{i}\bigg) = \sum_{i=1}^{k} m(J_{i}) = \sum_{i=1}^{k} \frac{1}{k} m(J) = m(J).
\end{displaymath}

Thus, we have proved that the two probability measures $m$ and $f_{\star}m$ coincide on the subintervals of $[0,1]$. In order to conclude that they are indeed the same, we invoke the following general result:

\begin{theorem}[Uniqueness of measures, Theorem 5.7 in~\cite{SchillingMeasuresIntegralsMart}]\label{thm:uniqueness-measures}
Let $\mu$ and $\nu$ two measures on the measurable space $(X,\mc{B})$. Let $\ms{C}\subset 2^{X}$ be a family that generates $\mc{B}$, \ie~$\mc{B}=\sigma(\ms{C})$, and suppose that
  \begin{itemize}
    \item if $A,B\in\mc{C}\implies A\cap B\in\ms{C}$;
    \item there exists ${(C_{n})}_{n\geq 1}\subset\ms{C}$ such that $C_{n}\uparrow X$ (\ie~$\bigcup_{n\geq 1} C_{n}=X$) with $\mu(C_{n})=\nu(C_{n})<\infty$, for every $n\geq 1$.
  \end{itemize}

  Then, $\mu=\nu$.
\end{theorem}

\subsection{Bernoulli shifts}
\label{sec:bernoulli-shifts}

Let $k$ be a natural number greater than $1$ and consider the space
\begin{displaymath}
  \Sigma_k^+:=\{1,\ldots,k\}^{\N_0}=\{(i_0i_1i_2\ldots) : i_n\in\{1,\ldots,k\},\ \forall n\geq 0\},
\end{displaymath}
endowed with the product topology considering the discrete topology on each factor.

Then the so called \emph{full shift} is the map $f\colon\Sigma_k^+\carr$ given by
\begin{displaymath}
  f(i_0i_1i_2\ldots):=(i_1i_2i_3\ldots), \quad\forall (i_n)\in\Sigma_k^+.
\end{displaymath}

\begin{exercise}\label{exer:full-shift-top-properties}
  Show that the full shift map $f$ is continuous and surjective and the set of periodic points
  \begin{displaymath}
    \Per(f):=\{x\in\Sigma_k^+ : \exists n\geq 1,\ f^n(x)=x\}
  \end{displaymath}
  is dense in $\Sigma_k^+$. On the other hand, there is $x_0\in\Sigma_k^+$ such that its orbit
  \begin{displaymath}
    \mc{O}_f^+(x_0):=\{f^n(x) : n\geq 0\}
  \end{displaymath}
  is dense in $\Sigma_k^+$.
\end{exercise}

As a consequence of Exercise~\ref{exer:full-shift-top-properties}, we see that there are infinitely many periodic orbits. However, there is another important family of invariant measures for full shifts, usually called \emph{Bernoulli measures.} In order to define them, let us considering an \textbf{stochastic vector} $\boldsymbol{p}=(p_1,p_2,\ldots,p_k)\in\R^k$, \ie~$p_i\geq 0$ and $\sum_i p_i=1$. Then, the \textbf{Bernoulli measure} induced by the stochastic vector $\boldsymbol{p}$ is the Borel probability measure $\mu$ on $\Sigma_k^+$ which is characterized by the following property:
\begin{displaymath}
  \mu[i_0,i_i,i_2,\ldots,i_m]:=p_{i_0}p_{i_1}\ldots p_{i_m}, \quad\forall (i_0,i_1,\ldots,i_m)\in\{1,\ldots,k\}^{m+1},
\end{displaymath}
and where $[i_0,i_1,\ldots,i_m]:=\big\{(j_0j_1j_2\ldots)\in\Sigma_k^+ : j_\ell=i_\ell,\ \forall \ell\{0,1,2,\ldots,m\}\big\}$.

\begin{exercise}
  Show that $\mu$ is $f$-invariant.
\end{exercise}

The space $\Sigma_k^+$ has interesting topological properties. A topological space $X$ is said to be \textbf{totally disconnected} when given any pair of different points $x,y\in X$, there is an open and closed subset $A\subset X$ such that $x\in A$ and $y\in X\setminus A$. On the other hand, a point $x\in X$ is said to be an \textbf{isolated point} when the set $\{x\}$ is a neighborhood of $x$.

A \textbf{Cantor space} is a compact metric space which is totally disconnected and has no isolated points.

\begin{exercise}[Construction of Cantor spaces]\label{exer:Cantor-space-construction}
  Let $(X_n)_{n\geq 1}$ be a sequence of finite sets with $\sharp X_n\geq 2$, for every $n\geq 1$. Then, if we endow each set $X_n$ with the discrete topology, the product space $\prod_{n\geq 1} X_n$ is a Cantor space.
\end{exercise}

\begin{exercise}[Uniqueness of Cantor spaces]\label{exer:Cantor-spaces-uniqueness}
  Show that, up to homeomorphisms, there is only one Cantor space.
\end{exercise}



\subsection{Circle rotations}\label{sec:circle-rotations}

Given any $\alpha\in\R$, one can define the map $T_{\alpha}\colon [0,1]\carr$ by
\begin{displaymath}
  T_{\alpha}(x):=x+\alpha - \floor{x+\alpha}, \quad\forall x\in [0,1].
\end{displaymath}

We claim that the Lebesgue measure $m$ on $[0,1]$ is again $T_{\alpha}$-invariant, for every $\alpha\in\R$. In fact, it is enough to consider the case where $\alpha\in (0,1)$ and then, one can easily show that given $[a,b]\subset [0,1]$, it holds
\begin{align*}
  T_{\alpha}^{-1}[a,b]&=[a-\alpha, b-\alpha],\quad \text{when } \alpha\leq a\leq b\leq 1;\\
  T_{\alpha}^{-1}[a,b]&=[a-\alpha+1,b-\alpha+1], \quad\text{when } 0\leq a\leq b<\alpha;\\
 T_{\alpha}^{-1}[a,b]&=[a-\alpha+1,1)\cap [0,b-\alpha], \quad\text{when } 0\leq a<\alpha\leq b\leq 1.
\end{align*}

Circle rotations are particular examples of translations on compact groups.

\subsection{Translations on locally compact groups}\label{sec:transl-locally-compact-groups}

A \textbf{topological group} is a ternary $(G,\mathscr{T},\cdot)$ where $G$ is a non-empty set, $\mathscr{T}$ is a Hausdorff topology on $G$ and $(G,\cdot)$ is group such that the map $G\times G\ni (g,h)\mapsto g\cdot h^{-1}\in G$ is a continuous map.

Analogously, a \textbf{Lie group} is a ternary $(G,\mathscr{C},\cdot)$ where $G$ is a non-empty set, $\mathscr{C}$ is a finite dimensional $C^{\infty}$ Hausdorff structure on $G$, $(G,\cdot)$ is a group and the map $G\times G\ni (g,h)\mapsto g\cdot h^{-1}\in G$ is $C^{\infty}$. Observe that any Lie group is a locally compact topological group.

Given topological group $G$ and any element $g\in G$, one can define the \textbf{left and right translations} by $L_{g}\colon G\ni h\mapsto gh\in G$ and $R_{g}\colon G\ni h\mapsto hg^{-1}\in G$, respectively. Notice that $L_{g}$ and $R_{g}$ are homeomorphism on $G$; and when $G$ is a Lie group, they are $C^{\infty}$ diffeomorphisms.

\begin{definition}[Haar measure]
  Given a locally compact group $G$, a \textbf{(left) Haar measure} is a Radon measure $\mu$ on $G$ (\ie~$\mu$ is Borel and $\mu(K)<\infty$, for every compact $K\subset G$) such that $\mu$ is $L_{g}$-invariant, for every $g\in G$, and there is a Borel set $E\subset G$ such that $\mu(E)>0$.
\end{definition}

\begin{theorem}[Existence and characterization of Haar measures]\label{thm:Haar-measures}
  If $G$ is locally compact group, then there exists a left Haar measure $\mu_{G}$. The support of $\mu_{G}$ is $G$ (\ie~$\mu_{G}(U)>0$, for every non-empty open subset $U\subset G$) and $\mu_{G}(G)<\infty$ if and only if $G$ is compact.

  On the other hand, if $\nu$ is another left Haar measure, then there exists $c>0$ such that
  \begin{displaymath}
    \mu(E)=c\nu(E), \quad\forall E\in\mc{B}(G).
  \end{displaymath}
\end{theorem}

\begin{proof}
  See for instance Halmos~\cite[\S58]{HalmosMeasureTheory}
\end{proof}

\begin{exercise}
  If $G$ is a Lie group, then show that any left Haar measure is a smooth measures (as defined in \S\ref{sec:smooth-measures}).
\end{exercise}

\begin{example}
  The space $\R^{d}$ is an abelian Lie group and $\Z^{d}<\R^{d}$ is a (normal) subgroup. Hence, $\T^{d}:=\R^{d}/\Z^{d}$ is a compact abelian Lie group, and there exists a unique left Haar measure $\lambda^{d}$ which is a probability measure. The measure $\lambda^{d}$ is left and right invariant is usually called the Haar or the Lebesgue measure on $\T^{d}$.
\end{example}

Given a topological group $G$, we say that a map $f\colon G\carr$ is a \textbf{topological group endomorphism} (or just an \textbf{endomorphism} for short) when $f$ is continuous and $f(xy)=f(x)f(y)$, for every $x,y\in G$. An endomorphism $f\colon G\carr$ is said to be a \textbf{topological group automorphism} (or just an \textbf{automorphism}) when $f$ is bijective and the map $f^{-1}$ is an endomorphism as well.

\begin{exercise}\label{exer:Haar-compact-groups}
  Let $G$ be a compact topological group, $\mu$ be the left Haar probability measure and $f\colon G\carr$ be an automorphism.
  \begin{enumerate}[(a)]
    \item Show that $\mu$ is a right Haar measure too, \ie~it is invariant by every right translation.
    \item Show that $\mu$ is an $f$-invariant measure.

    \item What happens when $f$ is just an endomorphism?

    \item And what happens when $G$ is non-compact and just locally compact?
  \end{enumerate}
\end{exercise}


\subsubsection{Compact abelian groups}
\label{sec:compact-abelian-groups}

Let $(G,+)$ be a compact abelian topological group. By Theorem~\ref{thm:Haar-measures}, we know that every Haar measure on $G$ is finite, by Exercise~\ref{exer:Haar-compact-groups}, every left Haar measure is a right one. So, there exists a unique Haar probability measure $\mu$ on $G$.

One can show that $G$ is metrizable. So, let $d'\colon G\times G\to \R_{0}^+$ be a distance function on $G$ which is compatible with the topology of $G$.  Then, we have

\begin{exercise}[Existence of invariant distance]\label{exer:abelian-groups-existence-invariant-distance}
  Show that the function $d\colon G\times G\to\R_0^+$ given by
  \begin{displaymath}
    d(g_1,g_2):=\int_G d'(g_1+h,g_2+h)\dd\mu(h), \quad\forall g_1,g_2\in G,
  \end{displaymath}
  is a distance function which is compatible with the topology of $G$ and $L_h$ is an isometry for $d$ for every $h\in G$ , \ie~it holds
  \begin{displaymath}
    d(g_1+h,g_2+h)=d(g_1,g_2), \quad\forall g_1,g_2,h\in G.
  \end{displaymath}
\end{exercise}

\begin{exercise}[Dynamics of translations]\label{exer:abelian-groups-dynamics-translations}
  For every $h\in G$, the following statements are equivalent:
  \begin{enumerate}[(i)]
    \item $L_h$ is \textbf{topologically transitive,} \ie~there is $g\in G$ such that its orbit $\mc{O}_{L_h}(g):=\{L_h^n(g) : n\in\Z\}$ is dense in $G$;
    \item $L_h$ is \textbf{minimal,} \ie~for every $g\in G$, the positive semi-orbit $\mc{O}_{L_h}^+(g):=\{L_h^n(g) : n\in\N_0\}$ is dense in $G$.
  \end{enumerate}
\end{exercise}


\subsubsection{Adding machines}
\label{sec:adding-machines}

There is a particularly interesting family of compact abelian groups with minimal translations, which are usually called \emph{adding machines.} To construct some of these spaces, let us start considering a number $p\in\N$, with $p\geq 2$, and the set
\begin{displaymath}
  \hat{\Z}_p:=\{0,1,\ldots,p-1\}^{\N_0},
\end{displaymath}
endowed with the product topology given by the discrete topology on each factor. By Exercise~\ref{exer:Cantor-space-construction}, we know that $\hat\Z_p$ is a Cantor space.
Then, consider the function $\phi\colon\N_0\to\hat\Z_p$ given by $\phi(n)=(a_0,a_1,a_2\ldots)\in\hat\Z_p$ if and only if there is $m\in\N_0$ such that $a_n=0$, for every $n\geq m$ and $n=\sum_{i=0}^\infty a_ip^i$. Writing $n$ in base $p$, one can easily check that function $\phi$ is indeed well defined, injective and its image is dense in $\hat\Z_p$. We can use this function $\phi$ to endow $\hat\Z_p$ with an abelian group structure, turning $\hat\Z_p$ into a compact abelian topological group:

\begin{exercise}\label{exer:p-adic-integer-group-structure-def}
  Show that there is a unique continuous function $+\colon\hat\Z_p\times\hat\Z_p\to\hat\Z_p$ such that
  \begin{displaymath}
    \phi(m)+\phi(n)=\phi(m+n), \quad\forall m,n\in\N_0.
  \end{displaymath}

  Then, show that $(\hat\Z_p,+)$ is indeed an abelian compact topological group, where $+$ is the function we defined above.
\end{exercise}

As a straight forward consequence of the denseness of the image of $\phi$ and Exercise~\ref{exer:p-adic-integer-group-structure-def}, we get that the translation $L_{\phi(1)}\colon\hat\Z_p\carr$ is transitive (because the semi-orbit $\mc{O}_{L_{\phi(1)}}^+\big(\phi(0)\big)=\phi(\N_0)$ is dense in $\hat\Z_p$), and hence, by Exercise~\ref{exer:abelian-groups-dynamics-translations}, we conclude that the translation $L_{\phi(1)}$ is minimal.

\subsection{The Gauss map}\label{sec:gauss-map}

Let us consider the map $G\colon (0,1)\to [0,1)$ given by
\begin{displaymath}
  G(x):=\frac{1}{x}-\floor{\frac{1}{x}}, \quad\forall x\in (0,1).
\end{displaymath}

\begin{exercise}
  Show that given $x\in (0,1)$ there exists $n\in\N$ such that $G^{n}(x)= 0$, if and only if $x\in\Q$.
\end{exercise}

So, if we define $X=(0,1)\setminus\Q$, then the map $G\big|_{X}\colon X\carr$ is surjective.

A remarkable fact is that there exists a smooth probability measure $\mu$ on $(0,1)$ (or on $X$) which is $G$-invariant. In fact we have

\begin{theorem}
  The measure $\mu$ given by
  \begin{displaymath}
    \mu(E):=\frac{1}{\log 2}\int_{E} \frac{\dd x}{1+x}, \quad\forall E\in\mc{B}(0,1),
  \end{displaymath}
  is a $G$-invariant Borel probability measure.
\end{theorem}

\begin{proof}
  It clearly holds $\mu(0,1)=1$. Let us prove that $\mu$ is $G$-invariant. For each $n\in\N$, consider the interval
  \begin{displaymath}
    I_{n}:=\bigg(\frac{1}{n+1},\frac{1}{n}\bigg), \quad\forall n\in\N.
  \end{displaymath}

  Observe that $G_{n}:=G\big|_{I_{n}}\colon I_{n}\to (0,1)$ is a $C^{\infty}$ diffeomorphism, for each $n\in\N$, and $m\Big(\bigsqcup_{n} I_{n})=m(0,1)=1$, where $m$ denotes the Lebesgue measure on $(0,1)$. Given any $\rho\in L^{1}((0,1),m)$, with $\rho\geq 0$ $m$-almost everywhere, let us consider the measure $\mu_{\rho}$ given by
  \begin{displaymath}
    \mu_{\rho}(E):=\int_{E} \rho\dd m, \quad\forall E\in \B(0,1).
  \end{displaymath}

  If we define the so called \emph{Perron-Frobenious operator} $\mathcal{L}$ by
  \begin{displaymath}
    \mathcal{L}\rho(y) :=\sum_{k=1}^{\infty}\frac{\rho(G_{k}^{-1}(y))}{G_{k}^{\prime}\big(G_{k}^{-1}(y)\big)}, \quad\forall y\in \bigsqcup_{k\geq 1} I_{k},
  \end{displaymath}
  we claim that
  \begin{displaymath}
    G_{\star}\mu_{\rho}(E) = \mu_{\mathcal{L}\rho}(E) = \int_{E} \mathcal{L}\rho\dd m, \quad\forall E\in\B(0,1).
  \end{displaymath}

  In order to prove this, let $E$ be an arbitrary Borel subset of $(0,1)$ and let us write $\ds{1}_{E}$ to denote its indicator function. By the classical change of variable formula we have~
  \begin{displaymath}
    \begin{split}
      \mu_{\rho}\big(G^{-1}(E)\big) &=\int_{0}^{1}\ds{1}_{G^{-1}(E)}(x)\rho(x)\dd m(x) = \int_{0}^{1}\ds{1}_{E}\big(G(x)\big)\rho(x)\dd m(x) \\
                                    &=\sum_{k=1}^{\infty}\int_{I_{n}} \ds{1}_{E}\big(G(x)\big)\rho(x)\dd m(x) \\
                                    &= \sum_{k=1}^{\infty}\int_{0}^{1}\ds{1}_{E}(y)\rho\big(G_{k}^{-1}(y)\big) {\big(G_{k}^{-1}\big)}^{\prime}(y) \dd m(y) \\
                                    &= \sum_{k=1}^{\infty}\int_{0}^{1} \ds{1}_{E}(y) \frac{\rho\big(G_{k}^{-1}(y)\big)}{G_{k}^{\prime}\big(G_{k}^{-1}(y)\big)}\dd m(y) \\
                                    &= \int_{0}^{1}\ds{1}_{E}(y) \mathcal{L}(\rho)(y)\dd m(y) = \mu_{\mathcal{L}\rho}(E),
    \end{split}
  \end{displaymath}
  where the last line of the chain of equalities follows by classical Levi's monotone convergence theorem.

  So, if we consider $\rho(x):=\frac{1}{1+x}$, the fact that $\mu_{\rho}$ is a $G$-invariant (finite) Borel measure follows from the fact that $\mathcal{L}\rho=\rho$ and this easily follows from a straightforward computation.
\end{proof}

\subsection{Some important exercises from~\cite[\S1.3.7]{VianaOliveiraFoundErgTh}}
\label{sec:some-import-exerc-1.3.7}

1.3.2, 1.3.4, 1.3.5, 1.3.6, 1.3.7, 1.3.9, 1.3.10, 1.3.11, 1.3.12.


\section{Constructing new measure preserving systems from old ones}\label{sec:constr-new-measure-pres-sys-old-ones}

\subsection{Product maps}\label{sec:product-maps}

Let $(X_{1},\B_{1},\mu_{1})$ and $(X_{2},\B_{2},\mu_{2})$ be two $\sigma$-finite measure spaces. The \textbf{product $\sigma$-algebra} $\B_{1}\otimes\B_{2}$ on $X_{1}\times X_{2}$ is defined as the $\sigma$-algebra generated by the family of \emph{rectangles} $\{E_{1}\times E_{2} : E_{i}\in\B_{i},\ i\in\{1,2\}\}$. Then, one shows that there exists a unique measure $\mu_{1}\otimes\mu_{2}$ on the measurable space $(X_{1}\times X_{2},\B_{1}\otimes\B_{2})$, called the \emph{product measure}, such that
\begin{displaymath}
  \mu_{1}\otimes\mu_{2}(E_{1}\times E_{2})=\mu_{1}(E_{1})\mu_{2}(E_{2}), \quad\forall E_{1}\in\B_{1},\ \forall E_{2}\in\B_{2}.
\end{displaymath}

When $f_{i}\colon (X_{1},\B_{i},\mu_{i})\carr$, for $i\in\{1,2\}$, are measure preserving maps, then invoking Theorem~\ref{thm:uniqueness-measures} one can easily show that the measure $\mu_{1}\otimes\mu_{2}$ is $f_{1}\times f_{2}$-invariant, where $f_{1}\times f_{2}\colon X_{1}\times X_{2}\ni (x_{1},x_{2})\mapsto \big(f_{1}(x_{1}),f_{2}(x_{2})\big)\in X_{1}\times X_{2}$.

By induction, this construction can be easily extended from two to any (finite) number of factor maps. Moreover, this construction can be extended to countably many factors, assuming that each factor is a probability space (see~\cite[Appendix A.2]{VianaOliveiraFoundErgTh} for details).


\subsection{Induction or renormalization}\label{sec:induct-or-renormalization}

Let $(X,\B,\mu)$ be a probability space, $f\colon X\carr$ be a measure preserving map and $A\in\B$ a measurable set with $\mu(A)>0$. By Poincaré recurrence theorem (Corollary~\ref{cor:return-infinitely-many-times-Poincare}) we know that $\mu$-almost every point $x\in A$ returns to $A$ infinitely many times. So we can define the function $\rho_{A}\colon A\to\N$ by
\begin{displaymath}
  \rho_{A}(x):=\min\left\{n\in\N : f^{n}(x)\in A\right\},
\end{displaymath}
and this function is well-defined $\mu$-a.e.~on $A$. Then we define the \textbf{first return map to $A$} by
\begin{displaymath}
  f_{A}(x):=f^{\rho_{A}(x)}(x),
\end{displaymath}
which is well define $\mu$-a.e.~on $A$. In order to see that $f_{A}$ is a well defined dynamical system, is enough to observe that if we define
\begin{displaymath}
  A_{f}:=A\cap \left(\bigcap_{m\geq 0} \bigcup_{n\geq m} f^{-n}(A)\right),
\end{displaymath}
then by Corollary~\ref{cor:return-infinitely-many-times-Poincare} one have that $\mu(A_{f})=\mu(A)>0$, with $A_{f}\subset A$, and one can easily check that $\rho_{A}(x)<\infty$  and $f_{A}(x)\in A_{f}$, for every $x\in A_{f}$. So, this means that $f_{A}\colon A_{f}\carr$ is a well defined dynamical system. Moreover, if define the probability measure $\mu_{A}$ on the measure space $(A,\mc{B}_{A})$ by
\begin{displaymath}
  \mu_{A}(E):=\frac{\mu(E)}{\mu(A)}, \quad\forall E\in\B_{A},
\end{displaymath}
and where $\B_{A}:=\{B\in\B : B\subset A_{f}\}$, we have the following

\begin{theorem}\label{thm:inductive-map-is-mpm}
  The map $f_{A}\colon A_{f}\carr$ is a $\B_{A}$-measurable and the measure $\mu_{A}$ is $f_{A}$-invariant.
\end{theorem}

\begin{proof}
  For the sake of simplicity of the notation, from now on we shall assume that $A=A_{f}$, and for each $n\in\N$, let us consider the set $A_{n}:=\rho^{-1}(n)\subset A$.
  Observe that
  \begin{displaymath}
    A=\bigsqcup_{n\in\N} A_{n},
  \end{displaymath}
  and for each $E\subset A$, it holds
  \begin{equation}
    \label{eq:f-A-inv-on-rho-level-sets}
    f_{A}^{-1}(E)=\bigsqcup_{n\in\N} f_{A}^{-1}(E)\cap A_{n} = \bigsqcup_{n\in\N} f^{-n}(E)\cap A_{n}.
  \end{equation}

  On the other hand, one can easily check that
  \begin{equation}
    \label{eq:rho-return-times-level-sets}
    A_{n}=\rho^{-1}(n)=A\cap \bigg(\bigcap_{i=1}^{n-1} f^{-i}(X\setminus A)\bigg) \cap f^{-n}(A), \quad\forall n\in\N.
  \end{equation}

  From~\eqref{eq:rho-return-times-level-sets} we conclude $\rho$ is a measurable function and then, by~\eqref{eq:f-A-inv-on-rho-level-sets} we get that $f_{A}$ is a $\B_{A}$-measurable map.

  Finally, in order to show that $\mu_{A}$ is $f_{A}$-invariant, it is enough to show that $\mu(E)=\mu(f^{-1}_{A}(E))$, for any $E\in\B_{A}$. In order to prove that, let us fix a measurable set $E\in\B_{A}$ and observe that, by~\eqref{eq:f-A-inv-on-rho-level-sets} we get
  \begin{displaymath}
    \mu\big(f_{A}^{-1}(E)\big) = \sum_{n=1}^{\infty}\mu\big(f^{-n}(E)\cap A_{n}\big).
  \end{displaymath}
  Let us prove that the right hand side of this equality is equal to $\mu(E)$. To do that, let us start noticing that
  \begin{equation}
    \label{eq:mu-E-equal-m-f-1-E-partition-A}
    \begin{split}
      \mu(E)=\mu\big(f^{-1}(E)\big) &= \mu\big(f^{-1}(E)\cap A\big) + \mu\big(f^{-1}(E)\setminus A\big) \\
      &= \mu\big(f^{-1}(E)\cap A_{1}\big) + \mu\big(f^{-1}(E)\setminus A\big),
    \end{split}
  \end{equation}
  where the first equality follows from the fact $f^{-1}(E)=\big(f^{-1}(E)\cap A\big)\sqcup \big(f^{-1}(E)\setminus A\big)$, and the second one from $f^{-1}(E)\cap A=f^{-1}(E)\cap A_{1}$.

  Then, focusing on the last term of~\eqref{eq:mu-E-equal-m-f-1-E-partition-A} we get
  \begin{equation}
    \label{eq:mu-E-f-1-E-setminus-A-second-argument}
    \begin{split}
      \mu\big(f^{-1}(E)\setminus A\big) &= \mu\big(f^{-2}(E)\setminus f^{-1}(A)\big) \\
                                        &= \mu\Big(\big(f^{-2}(E)\setminus f^{-1}(A)\big)\cap A \Big) + \mu\Big(\big(f^{-2}(E)\setminus f^{-1}(A))\setminus A\Big)\\
      &= \mu\big(f^{-2}(E)\cap A_{2}) + \mu\Big(f^{-2}(E)\setminus\big(A\cup f^{-1}(A)\big)\Big),
    \end{split}
  \end{equation}
  where the first equality follows from the invariance of $\mu$ by $f$, the second one by the partition argument we used above to prove~\eqref{eq:mu-E-equal-m-f-1-E-partition-A}, and the last equality follows from the fact that $\big(f^{-2}(E)\setminus f^{-1}(A)\big)\cap A = f^{-1}(E)\cap A_{2}$.

  Combining equations~\eqref{eq:mu-E-equal-m-f-1-E-partition-A} and~\eqref{eq:mu-E-f-1-E-setminus-A-second-argument} and repeating inductively the last argument, one can show that
  \begin{displaymath}
    \mu(E)= \sum_{i=1}^{n} \mu\big(f^{-i}(E)\cap A_{i}\big) + \mu\Big(f^{-n}(E)\setminus\big(A\cup f^{-1}(A)\cup\cdots\cup f^{-(n-1)}(A)\big)\Big),
  \end{displaymath}
  for every $n\geq 1$.

  So, in order to prove that $\mu(E)=\mu\big(f^{-1}(E)\big)$ it is enough to show that
  \begin{displaymath}
    \mu\bigg(f^{-n}(E) \setminus \Big(\bigcup_{i=1}^{n-1} f^{-i}(A) \Big)\bigg)\leq  \mu\bigg(f^{-n}(A) \setminus \Big(\bigcup_{i=1}^{n-1} f^{-i}(A) \Big)\bigg) \to 0
  \end{displaymath}
  as $n\to +\infty$. But this easily follows from the fact that $f^{-1}(A)\setminus A, f^{-2}(A)\setminus (A\cup f^{-1}(A))\ldots$ are pairwise disjoint and hence,
  \begin{displaymath}
    1=\mu(X) \geq \sum_{n=1}^{\infty}  \mu\bigg(f^{-n}(A) \setminus \Big(\bigcup_{i=1}^{n-1} f^{-i}(A) \Big)\bigg).
  \end{displaymath}
\end{proof}


\subsection{Kakutani-Rokhlin towers}\label{sec:Kakutani-Rokhlin-towers}

In this section we shall try to revert the previous construction we performed in \S\ref{sec:induct-or-renormalization}. To do that, for the time being let $(X,\B,\mu)$ be an arbitrary measure space, $\rho\colon X\to\N_{0}$ be an integrable function and $f\colon X\carr$ be a measure preserving map.

We first onsider the new measure space $(X_{\rho},\B_{\rho},\mu_{\rho})$ where each object is defined as follows: we first define the space $X_{\rho}$ by
\begin{displaymath}
  X_{\rho}:=\bigsqcup_{n\in\N_{0}} \bigsqcup_{0\leq i\leq n} \rho^{-1}(n)\times\{i\} \subset X\times\N_{0};
\end{displaymath}
the $\sigma$-algebra $\B_{\rho}$ is defined as the one generated by the family
\begin{displaymath}
  \B_{\rho}:=\sigma\left\{A\times\{i\}\subset X\times\N_{0} : A\in\B,\ \exists n\in\N_{0},\ 0\leq i\leq n,\ A\subset\rho^{-1}(n)\right\};
\end{displaymath}
and the measure $\mu_{\rho}$ is characterized as the only measure on the measurable space $(X_{\rho},\B_{\rho})$ such that
\begin{displaymath}
  \mu_{\rho}(A\times\{i\}):={\bigg(\int_{X}\rho+1\dd\mu\bigg)}^{-1}\mu(A),
\end{displaymath}
for every generating set $A\times\{i\}$ of $\B_{\rho}$.

\begin{exercise}
  Show that $(X_{\rho},\B_{\rho},\mu_{\rho})$ is a well-defined probability space.
\end{exercise}

Then we define the measure preserving map $f_{\rho}\colon X_{\rho}\carr$ by
\begin{displaymath}
  f_{\rho}(x,i):=
  \begin{cases}
    (x,i+1), & \text{for } x\in X,\ 0\leq i<\rho(x); \\
    \big(f(x),0\big), & \text{for } x\in X,\ i=\rho(x).
  \end{cases}
\end{displaymath}


\begin{exercise}
  Show that the map $f_{\rho}\colon (X_{\rho},\B_{\rho},\mu_{\rho})\carr$ is indeed a measure preserving map. Now consider the first return map (also called induction as in \S\ref{sec:induct-or-renormalization}) of the map $f_{\rho}$ to the set $A:=X\times\{0\}$. Show that $\mu_{\rho}(A)>0$ and the first return map is ``isomorphic'' to the original $f\colon (X,\B,\mu)\carr$. Can you give a definition of \emph{isomorphism} between measure preserving maps?
\end{exercise}

\subsection{Some important exercises from~\cite[\S1.4.4]{VianaOliveiraFoundErgTh}}
\label{sec:some-import-exerc-1.4.4}

1.4.3, 1.4.4.


\section{Existence of invariant measures for continuous maps}\label{sec:exist-invar-meas-cont-maps}

The main purpose of this section consists in proving the following

\begin{theorem}\label{thm:existence-inv-measures-cont-maps}
  Let $X$ be compact metric space and $f\colon X\carr$ be a continuous map. Then, there exists an $f$-invariant Borel probability measure.
\end{theorem}

In order to prove Theorem~\ref{thm:existence-inv-measures-cont-maps} we first need to introduce some concepts in Functional Analysis.

\subsection{Linear operators and dual spaces}\label{sec:line-oper-dual-spaces}

Let $(E,\norm{\cdot}_{E})$ and $(F,\norm{\cdot}_{F})$ be two Banach spaces over the field $\bb{K}=\R,\C$. A linear map $A\colon E\to F$ is continuous if and only its \textbf{norm}
\begin{equation}
  \label{eq:norm-operator}
  \norm{A}:=\sup_{x\in X\setminus\{0\}} \frac{\norm{Ax}_{F}}{\norm{x}_{E}} <\infty.
\end{equation}

The space of linear continuous maps from $E$ to $F$ is a denoted by $\mathcal{L}(E,F)$ and it turns to be a Banach space when its endowed with the norm given by~\eqref{eq:norm-operator}.

When $F$ is the field itself $\bb{K}$, a continuous linear operator $\eta\colon E\to\bb{K}$ is called a continuous linear functional. The \textbf{dual space} $E^{\star}$ is defined as the linear space $\mathcal{L}(E,\bb{K})$. Then we can define the norm $\norm{\cdot}_{E^{\star}}$ on $E^{\star}$ by
\begin{displaymath}
  \norm{\eta}_{E^{\star}}:=\sup\left\{\abs{\eta(x)} : x\in E,\ \norm{x}_{E}=1\right\}, \quad\forall \eta\in E^{\star}.
\end{displaymath}

The linear space $E^{\star}$ turns into a Banach space when it is endowed with the norm $\norm{\cdot}_{E^{\star}}$. The topology on $E^{\star}$ induced by the norm $\norm{\cdot}_{E^{\star}}$ is usually called the \textbf{strong topology}.

On the other hand, any continuous linear operator $A\colon E\to F$ naturally induces the so callled \textbf{adjoint operator} $A\colon F^{\star}\to E^{\star}$ given by
\begin{displaymath}
  A^{\star}\eta(x):=\eta\big(A(x)\big),\quad\forall x\in E,\ \forall\eta\in F^{\star}.
\end{displaymath}

\begin{exercise}
  Show that the adjoint operator is indeed a linear continuous operator and $\norm{A}\geq\norm{A^{\star}}$.
\end{exercise}



\subsection{Weak topologies}\label{sec:weak-topologies}

Given a topological space $(Y,\tau_{Y})$, a non-empty set $X$ and an arbitrary family of maps $\ms{F}={\{f_{i}\colon X\to Y\}}_{i\in I}$, we define the \textbf{weak topology} on $X$ induced by $\tau_{Y}$ and $\ms{F}$ as the weakest (\ie~the coarsest) topology $\tau$ on $X$ such that $f_{i}\colon (X,\tau)\to (Y,\tau_{Y})$ is continuous, for every $i\in I$.

Now, if $(E,\norm{\cdot}_{E})$ denotes a Banach space over $\mathbb{K}=\R,\C$, and $E^{\star}$ is its dual space, we can define the so called \textbf{weak topology} on $E$ as the weak topology induced by the family $E^{\star}$, \ie~it is the weakest topology on $E$ such that every $\eta\in E^{\star}$ is continuous on $E$. In general, the weak topology is weaker than the Banach topology induced by the norm $\norm{\cdot}_{E}$.

\begin{exercise}
  Show that if $E$ is a Banach space, then norm and the weak topologies coincide if and only if $E$ has finite dimension.
\end{exercise}

Now, we can consider the space $E^{\star\star}:={(E^{\star})}^{\star}$ which is the dual space of the Banach space $(E^{\star},\norm{\cdot}_{E^{\star}})$. Analogously, we can endow the space $E^{\star}$ with the \textbf{weak topology} which is the weakest topology such that every element of $E^{\star\star}$ is continuous.

However, there shall need to endow the space $E^{\star}$ with a weaker topology which is defined as follows. One can show that there is a natural injection of $E$ into $E^{\star\star}$, \ie~a linear continuous injective map $j\colon E\to E^{\star\star}$ which is given by
\begin{displaymath}
  j(x)(\eta):=\eta(x), \quad\forall x\in E,\ \forall\eta\in E^{\star}.
\end{displaymath}
Then, the so called \textbf{weak-$\star$ topology} is the weakest topology on $E^{\star}$ such that every linear functional in $j(E)$ is continuous. Given a sequence ${(\phi_{n})}_{n\geq 1}\subset E^{\star}$, we shall write
\begin{displaymath}
  \eta_{n}\rightharpoonup \eta\in E^{\star}, \quad\text{as } n\to \infty,
\end{displaymath}
when the sequence ${(\eta_{n})}_{n\geq 1}$ converges to $\eta$ in the weak-$\star$ topology.

\begin{exercise}
  Show that the weak and weak-$\star$ topologies coincide on $E^{\star}$ if and only if $E$ is reflexive space, \ie~ $j(E)=E^{\star\star}$. On the other hand, show that that the \emph{closed unit ball} on a Banach space $E$, which is given by $\{x\in E : \norm{x}_{E}\leq 1\}$ is compact if and only if $E$ has finite dimension.
\end{exercise}


We shall need the following
\begin{theorem}[Banach-Alaoglu theorem]\label{thm:Banach-Alaoglu}
  The closed unit ball $\{\eta\in E^{\star} : \norm{\eta}_{E^{\star}}\leq 1\}$ is a compact set when $E^{\star}$ is with its weak-$\star$ topology.
\end{theorem}

\begin{proof}
  See for instance~\cite[Theorem 8.10]{EinsiedlerWardFunctionalAnalysisBook}.
\end{proof}

\subsection{The space of continuous functions}\label{sec:space-cont-functions}

Let $(X,d)$ be a compact metric space and write $C^{0}(X)$ to denote the space of continuous real functions on $X$. Since $X$ is compact, any continuous function is bounded and hence, we can define the norm
\begin{displaymath}
  \norm{\phi}_{0}:=\sup_{x\in X} \abs{\phi(x)}, \quad\forall \phi\in C^{0}(X).
\end{displaymath}

\begin{exercise}
  Show that $C^{0}(X)$ is a Banach space when it is endowed with the norm $\norm{\cdot}$.
\end{exercise}

Given an arbitrary continuous map $f\colon X\carr$, we can consider the \textbf{induced pull-back linear map,} also called the \textbf{Koopman operator} $U_{f}\colon C^{0}(X)\carr$  given by
\begin{displaymath}
  U_{f}(\phi):=\phi\circ f, \quad\forall \phi\in C^{0}(X).
\end{displaymath}

\begin{exercise}
  Show that $U_{f}$ is a continuous linear operator and its norm
  \begin{displaymath}
    \norm{U_{f}}:=\sup\left\{\norm{U_{f}(\phi)}_{0} : \phi\in C^{0}(X), \norm{\phi}_{0}=1\right\}= 1.
  \end{displaymath}
  Moreover, $\norm{U_{f}\phi}=\norm{\phi}$, for every $\phi\in C^{0}(X)$ if and only if $f$ is surjective.
\end{exercise}

On the other hand, we have the following
\begin{theorem}[Riesz representation theorem]\label{thm:Riesz-representation-theorem-measures}
  If $\eta\colon C^{0}(X)\to\R$ is a continuous linear functional on the Banach space $(C^{0}(X),\norm{\cdot})$, then there exist two finite Borel measures $\mu^{-}$ and $\mu^{+}$ on $X$ such that
  \begin{displaymath}
    \eta(\phi)=\int_{X}\phi\dd\mu^{+} - \int_{X}\phi\dd\mu^{-}, \quad\forall\phi\in C^{0}(X).
  \end{displaymath}
\end{theorem}

\begin{proof}
  See for instance Einsiedler and Ward~\cite[Theorem 7.44]{EinsiedlerWardFunctionalAnalysisBook}.
\end{proof}

Then, let us write $\M(X)$ for the space of Borel probability measures on $X$. By Theorem~\ref{thm:Riesz-representation-theorem-measures} we can consider $\M(X)$ as a subset of unit ball of the dual space ${C^{0}(X)}^{\star}$. On the other hand, by Exercise~\ref{exer:pull-back-function-push-forward-measures} we know that the restriction of the adjoint operator of $U_{f}$ to the space $\M(X)$ coincide with the push-forward operator $f_{star}\colon\M(X)\carr$. In particular, we conclude $f_{\star}\colon\M(X)\carr$ is continuous when $\M(X)$ is endowed with the weak-$\star$ topology that inherits from the inclusion $\M(X)\subset {C^{0}(X)}^{\star}$.

\begin{exercise}\label{exer:compactness-MX}
  Show that the space of Borel probability measures $\M(X)$ is a closed subset of the closed unit ball $\bar B:=\left\{\xi\in {C^{0}(X)}^{\star} : \norm{\xi}\leq 1\right\}$ when it is endowed with the weak-$\star$ topology. Then, invoking Theorem~\ref{thm:Banach-Alaoglu}, conclude that $\M(X)$ is compact with the weak-$\star$ topology.
\end{exercise}

Now we can complete the
\begin{proof}[Proof of Theorem~\ref{thm:existence-inv-measures-cont-maps}]
  Let $\nu$ be any Borel probability measure on $X$. For each $n\in\N$, let us consider the measure $\mu_{n}\in\M(X)$ given by
  \begin{displaymath}
    \mu_{n}:=\frac{1}{n}\sum_{j=0}^{n-1}f_{\star}^{j}\nu.
  \end{displaymath}

  By Exercise~\ref{exer:compactness-MX} we know that $\M(X)$ is compact for the weak-$\star$ topology. So, there exists a subsequent ${(\mu_{n_{j}})}_{j\geq 1}$ and $\mu\in\M(X)$ such that $\mu_{n_{j}}\rightharpoonup\mu$, as $n_{j}\to\infty$. The, since $f_{\star}\colon\M(X)\carr$ is continuous, we get $f_{\star}\mu_{n_{j}}\rightharpoonup f_{\star}\mu$, as $n_{j}\to\infty$. So, in order to show that $\mu$ is $f$-invariant, it is enough to show that $f_{\star}\mu_{n}-\mu_{n}\rightharpoonup 0$, as $n\to\infty$. In order to do that, let $\phi\in C^{0}(X)$ be arbitrary and notice that
  \begin{displaymath}
    \int_{x}\phi\dd(f_{\star}\mu_{n}-\mu_{n}) = \frac{1}{n}\sum_{j=0}^{n-1}\left(\int_{X}\phi\dd f_{\star}^{j+1}\nu -\int_{X}\phi\dd f^{j}_{\star}\nu\right) = \frac{1}{n}\int_{X}\phi\circ f^{n}-\phi\dd\nu\to 0,
  \end{displaymath}
  as $n\to+\infty$.
\end{proof}

\begin{exercise}
  Let $X$ and $Y$ be two compact metric spaces, and $f\colon X\carr$ and $g\colon X\times Y\to Y$ two continuous functions. Let us consider the so called \emph{skew-product map} $H\colon X\times Y\carr$ given by
  \begin{displaymath}
    H(x,y):=\big(f(x),g(x,y)\big), \quad\forall (x,y)\in X\times Y.
  \end{displaymath}

  Given any invariant measure $\mu\in\M(f)$, let us show that there exists an invariant measure $\hat\mu\in\M(H)$ such that $\pi_{\star}\hat\mu=\mu$, where $\pi\colon X\times Y\to X$ denotes the projection on the first coordinate.
\end{exercise}


\subsection{Some important exercises from~\cite[\S2.1.6]{VianaOliveiraFoundErgTh}}
\label{sec:some-import-exerc-2.1.6}

2.1.2, 2.1.3, 2.1.4, 2.1.6.

\subsection{Some important exercises from~\cite[\S2.2.1]{VianaOliveiraFoundErgTh}}
\label{sec:some-import-exerc-2.2.1}

2.2.1, 2.2.2, 2.2.3, 2.2.4, 2.2.5.


\section{Ergodic theorems}\label{sec:ergodic-theorems}

Let $(X,\B,\mu)$ be a probability space and $f\colon X\carr$ a measure preserving endomorphism. Let $E\in\B$ be a subset such that $\mu(E)>0$. Then by Poincare's recurrence theorem (Corollary~\ref{cor:return-infinitely-many-times-Poincare}) we know that the set $\{i\in\N_{0} : f^{i}(x)\in E\}$ is infinite for $\mu$-almost every point $x\in E$. This can be considered as a qualitative information. In this section we shall focus on the problem of finding quantitative estimates for the frequency of sojourn
\begin{displaymath}
  \frac{1}{n}\sharp\{i\in\N_{0} : i<n,\ f^{i}(x)\in E\} = \frac{1}{n}\sum_{i=0}^{n-1}\ds{1}_{E}\circ f^{i}(x),
\end{displaymath}
as $n\to\infty$, and where $\ds{1}_{E}\colon X\to\{0,1\}$ denotes the indicator function of the set $E$. To do that, we shall analyze the asymptotic behavior of so called \textbf{Birkhoff averages} as the following one:
\begin{displaymath}
  \Bs_{f}^{n}\phi:=\frac{1}{n}\sum_{j=0}^{n-1}\phi\circ f^{j},
\end{displaymath}
for an arbitrary integrable function $\phi\colon X\to\R$.

The first result concerning the asymptotic behavior of Birkhoff averages is the celebrated

\begin{theorem}[von Neumann Ergodic Theorem]\label{thm:von-Neumann-ergodic-thm}
  Let $(x,\B,\mu)$ be a probability space and $f\colon X\carr$ be a measure preserving endomorphism. Then, for every $\phi\in L^{2}(X,\B,\mu)$, there is a $\tilde\phi\in L^{2}(X,\B,\mu)$ such that
  \begin{displaymath}
    \frac{1}{n}\sum_{j=0}^{n-1}\phi\circ f^{j}\to \tilde\phi, \quad\text{in } L^{2}(X,\B,\mu),
  \end{displaymath}
  as $n\to +\infty$. Moreover, it holds
  \begin{displaymath}
    \tilde\phi(x)=\tilde \phi\big(f(x)\big), \quad\text{for }\mu\text{-a.e. } x\in X.
  \end{displaymath}
\end{theorem}

In order to prove Theorem~\ref{thm:von-Neumann-ergodic-thm}, we first need to revise some concepts on Lebesgue and Hilbert spaces.

\subsection{Revision on Lebesgue spaces}\label{sec:revision-Lebesgue-spaces}

Let $(X,\mc{B},\mu)$ be an arbitrary measure space. For each real number $p$ with $1\leq p<\infty$, let us consider the space $L^{p}(X,\mc{B},\mu)$ of (equivalence classes of) measurable functions $\phi\colon X\to\C$ such that $\abs{\phi}^{p}$ is integrable, and where two functions are equivalent if and only if they coincide $\mu$-almost everywhere on $X$. Then one can show that the $L^{p}(X,\B,\mu)$ is a Banach space when it is endowed with the norm
\begin{displaymath}
  \norm{\phi}_{p}:= {\left(\int_{X}\abs{\phi}^{p}\dd\mu\right)}^{1/p},\quad\forall \phi\in L^{p}(X,\B,\mu).
\end{displaymath}

On the other hand, we define the space $L^{\infty}(X,B,\mu)$ as the space of (equivalence classes of) measurable functions $\phi\colon X\to\C$ such that there exists $C>0$ satisfying $\abs{\phi(x)}\leq C$, for $\mu$-almost every $x\in X$, and where two functions are equivalent when they coincide $\mu$-almost everywhere. The space $L^{\infty}(X,\B,\mu)$ is a Banach space when it is endowed with the norm
\begin{displaymath}
  \norm{\phi}_{\infty}:=\inf\left\{C>0 : \abs{\phi(x)}\leq C,\ \text{for }\mu\text{-a.e. } x\in X\right\}, \quad\forall \phi\in L^{\infty}(X,\B,\mu).
\end{displaymath}

\begin{exercise}
  If $(X,\B,\mu)$ is a probability space, $1\leq p<q\leq \infty$ and $\phi\in L^{q}(X,\B,\mu)$, then show that $\phi\in L^{p}(X,\B,\mu)$ and $\norm{\phi}_{p}\leq \norm{\phi}_{q}$. On the other hand, show that
  \begin{displaymath}
    \norm{\psi}_{\infty}=\lim_{p\to\infty}\norm{\psi}_{p}, \quad\forall \psi\in L^{\infty}(X,\B,\mu).
  \end{displaymath}
\end{exercise}

\begin{exercise}
  Consider the map $\langle\cdot,\cdot\rangle\colon L^{2}(X,\B,\mu)\times L^{2}(X,\B,\mu)\to\C$ given by
  \begin{displaymath}
    \langle\phi,\psi\rangle :=\int_{X}\phi\overline{\psi}\dd\mu,\quad\forall \phi,\psi\in L^{2}(X,\B,\mu).
  \end{displaymath}
  Show that $\langle\cdot,\cdot\rangle$ is well defined, is a scalar product and it turns $L^{2}(X,\B,\mu)$ into a Hilbert space.

  A Banach space is said to be \emph{Hibertlizable} when one can define a scalar product on it that turns it into a Hilbert space and whose topology coincides with the original one given by the Banach structure.

  Show that the Lebesgue space $L^{p}(X,\B,\mu)$ is Hilbertlizable if and only if $p=2$.
\end{exercise}

Given a probability space $(X,\B,\mu)$, a measure preserving map $f\colon X\carr$, and a measurable function $\phi\colon X\to\C$, consider the new measurable function $U_{f}\phi:=\phi\circ f$. The map $U_{f}$ is called the \textbf{Koopman operator}.
\begin{exercise}\label{exer:Koopman-operators-are-isometries}
  For every $p\in[1,\infty]$, show that $U_{f}\phi\in L^{p}(X,\B,\mu)$, whenever $\phi\in L^{p}(X,\B,\mu)$. Moreover, show that $U_{f}\colon L^{p}(X,\B,\mu)\carr$ is a linear isometry, \ie~$\norm{U_{f}\phi}_{p}=\norm{\phi}_{p}$, for every $\phi\in L ^{p}(X,\B,\mu)$ and every $p\in[1,\infty]$.
\end{exercise}

\subsection{Revision of Hilbert space theory}\label{sec:revis-Hilbert-space-theory}

Let $(\ms{H},\langle\cdot,\cdot\rangle)$ be a Hilbert over the field $\bb{K}=\R,\C$. We will write $E<\ms{H}$ when $E$ is a closed linear subspace. Given any non-empty subset $F\subset\ms{H}$, we define its \textbf{orthogonal complement} $F^{\perp}$ by
\begin{displaymath}
  F^{\perp}:=\{x\in\ms{H} : \langle x,y\rangle = 0,\ \forall y\in F\}.
\end{displaymath}
Observe that $F^{\perp}<\ms{H}$ and $F\cap F^{\perp}$ is either empty or the singleton $\{0\}$.

We will need the following
\begin{theorem}[Projection on convex subsets]\label{thm:projection-convex-closed-subsets-Hilbert}
  Let $K\subset\ms{H}$ be a non-empty convex closed set. Then there is a continuous map $P_{K}\colon\ms{H}\to K$ such that for every $x\in\ms{H}$ the vector $P_{K}(x)$ is the only one in $K$ such that
  \begin{displaymath}
    \norm{x-P_{k}(x)}=\inf\{\norm{x-y} : y\in K\}.
  \end{displaymath}

  Moreover, if $K$ is closed linear subspace, then $P_{K}$ is a linear continuous map such that
  \begin{displaymath}
    \langle x-P_{K}(x),x\rangle = 0, \quad\forall x\in\ms{H}.
  \end{displaymath}
  In particular, it holds $\ker P_{K}=K^{\perp}$ and $K\oplus K^{\perp}=\ms{H}$.
\end{theorem}

\begin{proof}
  See for instance~\cite[Theorem 3.13 and Corollary 3.17]{EinsiedlerWardFunctionalAnalysisBook}.
\end{proof}

Let $\ms{H}^{\star}$ denote its dual space, \ie~the space of continuous linear functionals on $\ms{H}$. We have the following

\begin{theorem}[Riesz representation theorem for Hilbert spaces]\label{thm:Riesz-representation-thm-Hilber-spaces}
  For every $\eta\in\ms{H}^{\star}$, there exists a unique $v_{\eta}\in\ms{H}$ such that
  \begin{displaymath}
    \eta(x)=\langle x,v_{\eta}\rangle, \quad\forall x\in\ms{H}.
  \end{displaymath}
\end{theorem}

\begin{proof}
  See for instance~\cite[Corollary 3.19]{EinsiedlerWardFunctionalAnalysisBook}.
\end{proof}

Theorem~\ref{thm:Riesz-representation-thm-Hilber-spaces} lets us redefine the concept of adjoint operator on Hilbert spaces. Given two Hilbert spaces $(\ms{H},\langle\cdot,\cdot\rangle_{\ms{H}})$ and $(\ms{J},\langle\cdot,\cdot\rangle_{\ms{J}})$ over $\{\}bb{K}=\R,\C$, and a continuous linear operator $A\in\mc{L}(\ms{H},\ms{J})$, we define its \textbf{adjoint operator} $A^{\star}\colon\ms{J}\to\ms{H}$ by the following property:
\begin{displaymath}
  \langle Ax,y\rangle = \langle x, A^{\star}y\rangle, \quad\forall x\in\ms{H},\ \forall y\in\ms{J}.
\end{displaymath}

\begin{exercise}
  Show that $A^{\star}\in\mc{L}(\ms{J},\ms{H})$, for every $A\in\mc{L}(\ms{H},\ms{J})$.
\end{exercise}

A linear operator $U\in\mc{L}(\ms{H},\ms{J})$ is said to be an \textbf{isometry} when
\begin{displaymath}
  \norm{Ux}_{\ms{J}}=\norm{x}_{\ms{H}}, \quad\forall x\in\ms{H}.
\end{displaymath}

\begin{exercise}\label{exer:isometries-Hilber-spaces}
  Let $U\colon\ms{H}\to\ms{J}$ be a continuous linear operator. Show that the following statements are equivalent:
  \begin{enumerate}
    \item $U$ is an isometry;
    \item $\langle Ux,Uy\rangle_{\ms{J}}=\langle x,y\rangle_{\ms{H}}$, for every $x,y\in\ms{H}$;
    \item $U^{\star}U=I$, where $I\in\mc{L}(\ms{H})$ denotes the identity map.
  \end{enumerate}
\end{exercise}

A surjective isometry $U\in\mc{L}(\ms{H})$ is called a \textbf{unitary operator}. One can easily check that $U$ is a unitary operator if and only if $U^{\star} U = UU^{\star}=I$.

Now we can prove the so called von Neumann's ergodic theorem for isometries on Hilbert spaces:

\begin{theorem}\label{thm:von-Neumann-ergodic-isometries-Hilbert-spaces}
  Let $\ms{H}$ be a Hilbert space over $\bb{K}=\R,\C$ and $U\colon\ms{H}\to\ms{H}$ be a linear isometry. Let $\ms{I}:=\ker(U-I)<\ms{H}$ and $P_{\ms{I}}\colon\ms{H}\to\ms{I}$ the projection map given by Theorem~\ref{thm:projection-convex-closed-subsets-Hilbert}. Then, it holds
  \begin{displaymath}
    \frac{1}{n}\sum_{i=0}^{n-1} U^{i}x\to P_{\ms{I}}(x), \quad\forall x\in\ms{H},
  \end{displaymath}
  as $n\to+\infty$.
\end{theorem}

\begin{proof}
  Let us start considering the closed linear space
  \begin{displaymath}
    E=\overline{(U-I)\ms{H}}.
  \end{displaymath}
  Let us show that $\ms{I}=E^{\perp}$. In fact, first observe that for every $x\in\ms{H}$ and any $y\in\ms{I}$ we have
  \begin{displaymath}
    \langle Ux-x,y\rangle = \langle Ux,y\rangle - \langle x,y\rangle = \langle Ux, Uy\rangle - \langle x,y\rangle =0,
  \end{displaymath}
  where the last equality follows from the fact that $U$ is an isometry and Exercise~\ref{exer:isometries-Hilber-spaces}. So, $\ms{I}\subset E^{\perp}$.

  On the other hand, if $y\in E^{\perp}$, then we have
  \begin{displaymath}
    0=\langle Ux-x,y\rangle = \langle Ux,y\rangle - \langle x,y\rangle = \langle x, U^{\star} y\rangle - \langle x,y\rangle = \langle x, U^{\star}y-y\rangle,
  \end{displaymath}
  for every $x\in\ms{H}$. So, by Theorem~\ref{thm:Riesz-representation-thm-Hilber-spaces} we have that $U^{\star}y-y=0$. Then we have
  \begin{displaymath}
    \begin{split}
      \norm{Uy-y}^{2} &= \langle Uy-y,Uy-y\rangle = \norm{Uy}^{2} - \langle Uy,y\rangle - \langle y,Uy\rangle + \norm{y}^{2} \\
                      &= \norm{y}^{2} - \langle y,U^{\star}y\rangle - \langle U^{\star} y,y\rangle - \norm{y}^{2} \\
      & = \norm{y}^{2} - \norm{y}^{2} - \norm{y}^{2} +\norm{y}^{2} = 0.
    \end{split}
  \end{displaymath}
  So, $y\in\ms{I}$ and we have proved that $\ms{I}=E^{\perp}$. Hence, it holds $\ms{H}=\ms{I}\oplus E$ and in order to conclude the proof it is enough to show that
  \begin{equation}
    \label{eq:zero-Birkhoff-averages-limit-for-almost-coboundaries}
    \frac{1}{n}\sum_{i=0}^{n-1} U^{i}x\to 0, \quad\text{as } n\to+\infty,
  \end{equation}
  for every $x\in E$. To do that, let $x$ denote an arbitrary point of $E$ and $\eps>0$ be any positive real number. So, there exists a $y\in\ms{H}$ such that
  \begin{displaymath}
    \norm{x- (Uy-y)}<\frac{\eps}{2},
  \end{displaymath}
  and we get
  \begin{displaymath}
    \begin{split}
      \norm{\frac{1}{n}\sum_{i=0}^{n-1} U^{i}x} & = \frac{1}{n} \norm{\sum_{i=0}^{n-1} U^{i}\Big(\big(Uy-y\big) + x - (Uy-y)\Big)} \\
                                                &\leq \frac{1}{n} \norm{\sum_{i=0}^{n-1} U^{i+1}y-U^{i}y} + \frac{1}{n} \sum_{i=0}^{n-1} \norm{U^{i}\big(x - (Uy-y)\big)} \\
                                                &= \frac{\norm{U^{n}y-y}}{n} +  \frac{1}{n} \sum_{i=0}^{n-1} \norm{x - (Uy-y)} \leq \frac{2}{n}\norm{y} + \frac{\eps}{2}\leq \eps,
    \end{split}
  \end{displaymath}
  for every $n\geq \dfrac{4\norm{y}}{\eps}$. Since $\eps>0$ is arbitrary, we conclude that~\eqref{eq:zero-Birkhoff-averages-limit-for-almost-coboundaries} holds.
\end{proof}

\begin{proof}[Proof of Theorem~\ref{thm:von-Neumann-ergodic-thm}]
  It is a straightforward consequence of Theorem~\ref{thm:von-Neumann-ergodic-isometries-Hilbert-spaces} and the fact that the Koopman operator $U_{f}\colon L^{2}(X,\B,\mu)\carr$ is an isometry (Exercise~\ref{exer:Koopman-operators-are-isometries}).
\end{proof}

Here we review some interesting consequences of von Neumann's ergodic theorem (Theorem~\ref{thm:von-Neumann-ergodic-thm}):

\begin{corollary}\label{cor:von-Neumann-for-automorphisms-back-eq-for}
  Let $(X,\B,\mu)$ be a probability space and $f\colon X\carr$ be a measure preserving automorphism. Then, for every $\phi\in L^{2}(X,\B,\mu)$, there exists $\tilde\phi\in L^{2}(X,\B,\mu)$ such that
  \begin{displaymath}
    \lim_{n\to+\infty}\frac{1}{n}\sum_{j=0}^{n-1}\phi\circ f^{i} = \lim_{n\to-\infty}\frac{1}{n}\sum_{i=0}^{n-1}\phi\circ f^{-i} =\tilde\phi,
  \end{displaymath}
  where both limits are taken in $L^{2}$.
\end{corollary}

\begin{proof}
  This is a straightforward consequence of Theorem~\ref{thm:von-Neumann-ergodic-isometries-Hilbert-spaces} and the simple fact that $U_{f}$ and $U_{f^{-1}}={U_{f}}^{-1}$ are isometries, and $\ker(U_{f}-I)=\ker(U_{f^{-1}}-I)$
\end{proof}

\begin{corollary}\label{cor:von-Neumann-convergence-in-L1}
  Let $(X,\B,\mu)$ be a probability space, $f\colon X\carr$ be measure preserving map and $\phi\in L^{1}(X,\B,\mu)$. Then, there is $\tilde\phi\in L^{1}(X,\B,\mu)$ such that
  \begin{displaymath}
    \frac{1}{n}\sum_{i=0}^{n-1}\phi\circ f^{i}\to \tilde\phi, \quad\text{in }L^{1},
  \end{displaymath}
  as $n\to+\infty$.
\end{corollary}

\begin{proof}
  For the sake of simplicity of notation, given any $\psi\colon X\to\R$ and any $n\in\N$, let us write $\Bs_{f}^{n}\psi$ for the Birkhoff average of order $n$ of $\psi$ over $f$, \ie~we have
   \begin{displaymath}
    \Bs_{f}^{n}\psi = \frac{1}{n}\sum_{j=0}^{n-1}\psi\circ f^{j}.
  \end{displaymath}

  Let us fix a real number $\eps>0$. Since $L^2(X,\B,\mu)$ is dense in $L^1(X,\B,\mu)$, there is $\xi\in L^2(X,\B,\mu)$ such that
  \begin{displaymath}
    \norm{\phi-\xi}_1<\frac{\eps}{4},
  \end{displaymath}
  and then,
  \begin{displaymath}
    \norm{\Bs_f^n\phi-\Bs_f^n\xi}_1<\frac{\eps}{4}, \quad\forall n\geq 1.
  \end{displaymath}

  On the other hand, by Theorem~\ref{thm:von-Neumann-ergodic-thm} we know that there is $\tilde\xi\in L^{2}(X,\B,\mu)$ such that
  \begin{displaymath}
    \norm{\Bs_f^n\xi-\tilde\xi}_2\to 0, \quad\text{as } n\to+\infty,
  \end{displaymath}
  and thus, there is $N\in\N$ such that
  \begin{displaymath}
    \norm{\Bs_f^n\xi-\tilde\xi}_1\leq \norm{\Bs_f^n\xi-\tilde\xi}_2<\eps/4, \quad\forall n\geq N.
  \end{displaymath}
  Putting together all these estimates we get
  \begin{displaymath}
    \norm{\Bs_f^m\phi-\Bs_f^n\phi}_1\leq \norm{\Bs_f^m\phi-\Bs_f^m\xi}_1 + \norm{\Bs_f^m\xi-\tilde\xi}_1 + \norm{\tilde\xi - \Bs_f^n\xi}_1 + \norm{\Bs_f^n\xi-\Bs_f^n\phi}_1< \eps,
  \end{displaymath}
 for every $m,n\geq N$. So, ${\big(\Bs_f^n\phi\big)}_{n\geq 1}$ is a Cauchy sequence in $L^1(X,\B,\mu)$ and hence, there exists its limit $\tilde\phi\in L^1(X,\B,\mu)$.
\end{proof}

\subsection{Some important exercises from~\cite[\S3.1.4]{VianaOliveiraFoundErgTh}}
\label{sec:some-import-exerc-3.1.4}

3.1.1, 3.1.2, 3.1.3.


\subsection{Birkhoff ergodic theorem}
\label{sec:birkh-ergod-thm}

Now we are going to prove that indeed pointwise convergence holds on Corollary~\ref{cor:von-Neumann-convergence-in-L1}:

\begin{theorem}[Birkhoff ergodic theorem]\label{thm:Birkhoff-ergodic-thm}
  Let $(X,\B,\mu)$ be a probability space, $f\colon X\carr$ a measure preserving map and $\phi\colon X\to\R$ any integrable function. Then there is $\tilde\phi\in L^{1}(X,\B,\mu)$ such that
  \begin{displaymath}
    \Bs_{f}^{n}\phi(x)=\frac{1}{n}\sum_{j=0}^{n-1}\phi\big(f^{j}(x)\big)\to\tilde\phi(x), \quad\text{as } n\to\infty,
  \end{displaymath}
  for $\mu$-almost every $x\in X$.

  Moreover, it holds
  \begin{displaymath}
    \int_{X}\phi\dd\mu = \int_{X}\tilde\phi\dd\mu,
  \end{displaymath}
  and
  \begin{displaymath}
    \tilde\phi(x)=\tilde\phi\big(f(x)\big), \quad\text{for }\mu\text{-a.e. } x\in X.
  \end{displaymath}
\end{theorem}

In order to prove Theorem~\ref{thm:Birkhoff-ergodic-thm} we first need some preliminary results:

\begin{lemma}[Borel-Cantelli]\label{lem:Borel-Cantelli}
  Let $(X,\B,\mu)$ be a probability space and $A_{1},A_{2},\ldots\in\B$ be a sequence of measurable sets. If
  \begin{displaymath}
    \sum_{n\geq 1}\mu(A_{n})<\infty,
  \end{displaymath}
  then it holds
  \begin{displaymath}
    \mu\Big(\limsup_{n}A_{n}\Big)=0,
  \end{displaymath}
  where the \emph{limit superior} of the sequence $(A_{n})$ is defined by
  \begin{displaymath}
    \limsup_{n} A_{n} := \bigcap_{m\geq 1}\bigcup_{n\geq m} A_{n}.
  \end{displaymath}
\end{lemma}

\begin{proof}
  For each $m\in\N$, let us consider the set
  \begin{displaymath}
    B_{m}:=\bigcup_{n\geq m} A_{n}.
  \end{displaymath}
  Observe that $B_{1}\supset B_{2}\supset\cdots$ and since $\sum_{n\geq 1}\mu(A_{n})<\infty$, it holds $\mu(B_{m})\leq \sum_{n\geq m}\mu(A_{n})\to 0$, as $m\to\infty$. Finally, notice that
  \begin{displaymath}
    \mu\Big(\limsup_{n}A_{n}\Big) = \mu\bigg(\bigcap_{m\geq 1} B_{m}\bigg) = \lim_{m\to\infty} \mu(B_{m})=0.
  \end{displaymath}
\end{proof}

\begin{lemma}\label{lem:phi-fn-div-n-to-zero}
  Let $(X,\B,\mu)$ be a probability space and $\phi\in L^{1}(X,\B,\mu)$. Then it holds
  \begin{displaymath}
    \lim_{n\to+\infty}\frac{\phi\big(f^{n}(x)\big)}{n}=0, \quad\text{for }\mu\text{-a.e. }x\in X.
  \end{displaymath}
\end{lemma}

\begin{proof}
  For the sake of simplicity of notation, given any $\psi\colon X\to \R$ and any $A\subset\R$ and any $r\in\R$ let us write
  \begin{displaymath}
    [\psi \in A]:=\psi^{-1}(A), \quad\text{and}\quad [\psi>r]:=[\psi\in (r,\infty)].
  \end{displaymath}

  Then, first observe that given an arbitrary point $x\in X$, the sequence $\phi(f^{n}(x))/n$ does not converge to $0$ if and only if the limit superior of $\abs{\phi(f^{n}(x))}/n$ is positive. So, we have
  \begin{displaymath}
    \left\{x\in X : \frac{\phi\big(f^{n}(x)\big)}{n}\not\to 0,\ \text{as } n\to+\infty\right\} = \bigcup_{p\geq 1} \bigcap_{m\geq 1} \bigcup_{n\geq m} \left[\frac{\abs{\phi\circ f^{n}}}{n} > \frac{1}{p}\right].
  \end{displaymath}

  Thus in order to prove this lemma, it is enough to show that given any $\eps>0$,
  \begin{displaymath}
    \mu\left(\limsup_{n} \left[\frac{\abs{\phi\circ f^{n}}}{n} > \eps\right]\right) = 0,
  \end{displaymath}
  and by Lemma~\ref{lem:Borel-Cantelli}, to prove this last equality it is enough to show that
  \begin{displaymath}
    \sum_{n\geq 1} \mu\left[\frac{\abs{\phi\circ f^{n}}}{n} > \eps\right]<\infty.
  \end{displaymath}
  Finally, to prove the convergence of this series, observe that
  \begin{displaymath}
    \begin{split}
     \sum_{n\geq 1} \mu\left[\frac{\abs{\phi\circ f^{n}}}{n} > \eps\right] &= \sum_{n\geq 1} \mu\left(f^{-n}\left[\frac{\abs{\phi}}{n} >\eps\right]\right) = \sum_{n\geq 1} \mu\left[\frac{\abs{\phi}}{n}>\eps\right] \\
                                                                           & \sum_{n\geq 1} \mu\left[\frac{\abs{\phi}}{\eps}> n\right] = \sum_{n\geq 1} \sum_{j\geq n} \mu\left[j<\frac{\abs{\phi}}{\eps}\leq j+1\right] \\
                                                                           & = \sum_{j\geq 1} j\mu\left[j<\frac{\abs{\phi}}{\eps}\leq j+1\right] \leq \int_{X}\frac{\abs{\phi}}{\eps}\dd\mu<\infty.
    \end{split}
  \end{displaymath}
\end{proof}

\begin{theorem}[Maximal ergodic theorem]\label{thm:maximal-ergodic-thm}
  Let $(X,\B,\mu)$, $f$ and $\phi$ as in Theorem~\ref{thm:Birkhoff-ergodic-thm}. For any $\alpha>0$, let us consider the set
  \begin{displaymath}
    E^{\alpha}(\phi) := \left[\sup_{n\geq 1} \Bs_{f}^{n}\phi >\alpha\right] = \left\{x\in X : \sup_{n\geq 1} \frac{1}{n}\sum_{j=0}^{n-1}\phi\big(f^{j}(x)\big)>\alpha\right\}.
  \end{displaymath}
  Then it holds
  \begin{displaymath}
    \int_{E^{\alpha}(\phi)}\phi\dd\mu \geq \alpha\mu\big(E^{\alpha}(\phi)\big).
  \end{displaymath}
\end{theorem}

\begin{proof}[Proof of Theorem~\ref{thm:maximal-ergodic-thm}]
  Let us start noticing that given any $\alpha\in\R$ and any $\phi\in L^{1}(X,\B,\mu)$, it holds
  \begin{displaymath}
    E^{\alpha}(\phi) = E^{0}(\phi-\alpha),
  \end{displaymath}
  and therefore,
  \begin{displaymath}
    \int_{E^{0}(\phi-\alpha)}\phi-\alpha\dd\mu \geq 0\quad \iff\quad \int_{E^{\alpha}(\phi)}\phi\dd\mu\geq \alpha\mu(E^{\alpha}(\phi)).
  \end{displaymath}

  So, it is enough to prove this theorem for $\alpha=0$.

  For each $n\geq 1$, let us consider the function
  \begin{displaymath}
    \phi_n(x):=\sup\left\{\sum_{j=0}^{k-1}\phi\big(f^j(x)\big) : 1\leq k\leq n\right\}, \quad\forall x\in X.
  \end{displaymath}

  If we write $[\phi_n>0]:=\{x\in X : \phi_n(x)>0\}$, so it holds $[\phi_n>0]\subset[\phi_{n+1}>0]$, for every $n\geq 1$ and
  \begin{displaymath}
    E^0(\phi):=\bigcup_{n\geq 1}[\phi_n>0].
  \end{displaymath}

  Thus, to prove that $\int_{E^{0}(\phi)}\phi\dd\mu\geq 0$, by the dominated convergence theorem it is enough to show that
  \begin{displaymath}
    \int_{[\phi_n>0]}\phi\dd\mu \geq 0, \quad\forall n\geq 1.
  \end{displaymath}

  Then, given any $n\geq 1$ and any $m\in\{1,2,\ldots,n\}$, we know that
  \begin{displaymath}
    \phi_n(x)\geq \sum_{j=0}^{m-1}\phi\big(f^j(x)\big), \quad\forall x\in X,
  \end{displaymath}
  and hence,
  \begin{equation}
    \label{eq:phi-n-f-phi-geq-phi-n+1}
    \phi_n\big(f(x)\big) + \phi(x) \geq \sum_{j=0}^m\phi\big(f^j(x)\big), \quad\forall x\in X,\ \forall m\in\{1,\ldots,n\}.
  \end{equation}

  Fixing a point $x\in X$, let us consider the two possible cases: either $\phi_n\big(f(x)\big)\geq 0$ or $\phi_n\big(f(x)\big)< 0$.

  In the first case, when $\phi_n\big(f(x)\big)\geq 0$, we have $\phi_n\big(f(x)\big)+\phi(x)\geq \phi(x)$. Then, by estimate~\eqref{eq:phi-n-f-phi-geq-phi-n+1} it follows that
  \begin{displaymath}
    \phi_n\big(f(x)\big)+ \phi(x)\geq \phi_n(x), \quad\forall x\in [\phi_n\circ f\geq 0].
  \end{displaymath}

  In the second case, when $\phi_n\big(f(x)\big)<0$, by~\eqref{eq:phi-n-f-phi-geq-phi-n+1} we get
  \begin{displaymath}
    \phi(x)\geq \sum_{j=0}^m\phi\big(f^j(x)\big) - \phi_n\big(f(x)\big) > \sum_{j=0}^m\phi\big(f^j(x)\big), \quad\forall m\in\{1,\ldots,n\},
  \end{displaymath}
  and then, $\phi_n(x)=\phi(x)$, whenever $x\in [\phi_n\circ f<0]$.

  Putting together these estimates we get
  \begin{displaymath}
    \begin{split}
      \int_{[\phi_n>0]}\phi\dd\mu &= \int_{[\phi_n>0]\cap [\phi_n\circ f<0]}\phi\dd\mu + \int_{[\phi_n>0]\cap [\phi_n\circ f\geq 0]}\phi\dd\mu \\
      &\geq \int_{[\phi_n>0]\cap [\phi_n\circ f<0]}\phi\dd\mu + \int_{[\phi_n>0]\cap [\phi_n\circ f\geq 0]}\phi_n-\phi_n\circ f\dd\mu \\
      &= \int_{[\phi_n>0]\cap [\phi_n\circ f<0]}\phi_n\dd\mu + \int_{[\phi_n>0]\cap [\phi_n\circ f\geq 0]}\phi_n-\phi_n\circ f\dd\mu \\
      &= \int_{[\phi_n>0]}\phi_n\dd\mu - \int_{[\phi_n>0]\cap [\phi_n\circ f\geq 0]}\phi_n\circ f\dd\mu \\
      &\geq \int_{[\phi_n>0]}\phi_n\dd\mu - \int_{[\phi_n\circ f\geq 0]}\phi_n\circ f\dd\mu \\
      &= \int_{[\phi_n>0]}\phi_n\dd\mu - \int_{[\phi_n\geq 0]}\phi_n\dd(f_\star\mu) = 0.
    \end{split}
  \end{displaymath}
\end{proof}

\begin{corollary}\label{cor:maxima-ergodic-thm-on-invariant-set}
  If $(X,\B,\mu)$, $f$, $\alpha$ and $\phi$ are as in Theorem~\ref{thm:maximal-ergodic-thm} and $A\in\B$ satisfies $A\subset E^{\alpha}(\phi)$ and $f^{-1}(A)=A$, then it holds
  \begin{displaymath}
    \int_{A}\phi\dd\mu \geq \alpha\mu(A).
  \end{displaymath}
\end{corollary}

\begin{exercise}
  Prove Corollary~\ref{cor:maxima-ergodic-thm-on-invariant-set}.
\end{exercise}

Now we are ready for
\begin{proof}[Proof of Theorem~\ref{thm:Birkhoff-ergodic-thm}]
  Let us start considering the functions $\overline{\phi},\underline{\phi}\colon X\to[-\infty,+\infty]$ given by
  \begin{displaymath}
    \overline{\phi}(x):=\limsup_{n\to+\infty}\Bs_f^n\phi(x) \quad\text{and}\quad \underline{\phi}(x):=\liminf_{n\to+\infty} \Bs_f^n\phi(x),
  \end{displaymath}
  for every $x\in X$.

  \begin{exercise}
    Using Lemma~\ref{lem:phi-fn-div-n-to-zero}, show that the functions $\overline\phi$ and $\underline{\phi}$ are (essentially) $f$-invariant.
  \end{exercise}

  Our purpose now consists in showing that the functions $\overline{\phi}$ and $\underline{\phi}$ coincide $\mu$-a.e. In order to prove that, given real numbers $\alpha<\beta$, consider the set
  \begin{displaymath}
    E_\alpha^\beta:=\left\{x\in X : \underline{\phi}(x)<\alpha<\beta<\overline{\phi}(x)\right\},
  \end{displaymath}
  and observe that
  \begin{displaymath}
    [\underline{\phi}\neq \overline{\phi}] = [\underline{\phi}<\overline{\phi}] = \bigcup_{\alpha,\beta\in\Q,\ \alpha<\beta} E_{\alpha}^{\beta}
  \end{displaymath}
  So, to prove that $\underline{\phi}$ and $\overline{\phi}$ coincide $\mu$-a.e.~it is enough to show that $\mu\big(E_\alpha^\beta\big)=0$, for every pair of real numbers $\alpha<\beta$.

  Since $\underline{\phi}$ and $\overline{\phi}$ are $f$-invariant, it is easy to show that $f^{-1}\Big(E_\alpha^\beta\Big)=E_\alpha^\beta$ up to sets of $\mu$-measure zero. Then, observing that $E_\alpha^\beta\subset E^\beta(\phi)$, where the last set is defined as in Theorem~\ref{thm:maximal-ergodic-thm}, applying Corollary~\ref{cor:maxima-ergodic-thm-on-invariant-set} we can conclude
  \begin{displaymath}
    \int_{E_\alpha^\beta}\phi\dd\mu\geq\beta\mu\Big(E_\alpha^\beta\Big).
  \end{displaymath}


  On the other hand, notice that $E_{\alpha}^{\beta}\subset E^{-\alpha}(-\phi)$. So, we can again apply Corollary~\ref{cor:maxima-ergodic-thm-on-invariant-set} to the function $-\phi$ to get
  \begin{displaymath}
    \int_{E_\alpha^\beta}-\phi\dd\mu\geq-\alpha\mu\Big(E_\alpha^\beta\Big).
  \end{displaymath}
  Putting together these estimates we conclude
  \begin{displaymath}
    \beta\mu\Big(E_\alpha^\beta\Big) \leq \alpha\mu\Big(E_\alpha^\beta\Big),
  \end{displaymath}
  and since $\alpha<\beta$, this can only happend when $\mu\Big(E_\alpha^\beta\Big)=0$.

  The fact that $\tilde\phi:=\underline{\phi}=\overline{\phi}$ belongs to $L^{1}(X,\B,\mu)$ and
  \begin{displaymath}
    \int_{X}\phi\dd\mu=\int_{X}\tilde\phi\dd\mu,
  \end{displaymath}
  will be proved in Corollary~\ref{cor:Birkhoff-Lp-convergence}
\end{proof}

There are some important consequences of Birkhoff ergodic theorem:

\begin{corollary}\label{cor:Birkhoff-tilde-phi-in-Lp-when-phi-in-Lp}
  If $(X,\B,\mu)$ and $f$ are as in Theorem~\ref{thm:Birkhoff-ergodic-thm} and $\phi\in L^{p}(X,\B,\mu)$ with $1\leq p<\infty$, then the limit function $\tilde\phi$ given by Theorem~\ref{thm:Birkhoff-ergodic-thm} belongs to $L^{p}(X,\B,\mu)$ and it holds
  \begin{displaymath}
    \norm{\tilde\phi}_{p}\leq \norm{\phi}_{p}.
  \end{displaymath}
\end{corollary}


\begin{proof}
  Notice that
  \begin{displaymath}
    \abs{\tilde\phi(x)}\leq \lim_{n\to+\infty}\frac{1}{n}\sum_{j=0}^{n-1}\abs{\phi\circ f^{j}(x)}, \quad\text{for }\mu\text{-a.e. } x\in X,
  \end{displaymath}
  and hence,
  \begin{displaymath}
    \abs{\tilde\phi(x)}^{p}\leq \lim_{n\to+\infty}{\left(\frac{1}{n}\sum_{j=0}^{n-1}\abs{\phi\circ f^{j}(x)}\right)}^{p}.
  \end{displaymath}
  Then, as a consequence of Fatou lemma and Minkowski inequality we get
  \begin{displaymath}
    \begin{split}
      \norm{\tilde\phi}_{p}={\left[\int_{X}\abs{\tilde\phi(x)}^{p}\dd\mu\right]}^{\frac{1}{p}} &\leq \liminf_{n} {\left[\int_{X}{\left(\frac{1}{n}\sum_{j=0}^{n-1}\abs{\phi\big(f^{j}(x)\big)}\right)}^{p}\dd\mu\right]}^{\frac{1}{p}}\\
      & \leq \liminf_{n} \frac{1}{n}\sum_{j=0}^{n-1} {\left(\int_{X}\abs{\phi\big(f^{j}(x)\big)}^{p}\dd\mu\right)}^{\frac{1}{p}} =\norm{\phi}_{p}.
    \end{split}
  \end{displaymath}
\end{proof}

\begin{corollary}[$L^{p}$-convergence of ergodic averages]\label{cor:Birkhoff-Lp-convergence}
  If $(X,\B,\mu)$,  $f$, $\phi$ and $\tilde\phi$ are as in Corollary~\ref{cor:Birkhoff-tilde-phi-in-Lp-when-phi-in-Lp}, then
  \begin{displaymath}
    \norm{\Bs_{f}^{n}\phi-\tilde\phi}_{p}\to 0,
  \end{displaymath}
  as $n\to+\infty$,and where  $\Bs_{f}^{n}\phi=\dfrac{1}{n}\sum_{j=0}^{n-1}\phi\circ f^{j}$. Therefore, it holds
  \begin{displaymath}
    \int_{X}\phi\dd\mu=\int_{X}\tilde\phi\dd\mu.
  \end{displaymath}
\end{corollary}

\begin{proof}
  By Corollary~\ref{cor:Birkhoff-Lp-convergence} we know that $\tilde\phi\in L^{p}(X,\B,\mu)$.
  Then to prove the converges in $L^{p}$, let us start assuming $\phi\in L^{\infty}(x,\B,\mu)$. So, $\norm{\Bs_{f}^{n}\phi}_{\infty}\leq\norm{\phi}_{\infty}$, for every $n\geq 1$, and hence, $\tilde\phi\in L^{\infty}(X,\B,\mu)$, as well. So,
  \begin{displaymath}
    \norm{\Bs_{f}^{n}\phi-\tilde\phi}_{p}\leq \norm{\Bs_{f}^{n}\phi-\tilde\phi}_{\infty} \leq\norm{\phi}_{\infty}+\norm{\tilde\phi}_{\infty},
  \end{displaymath}
  and $\Bs_{f}^{n}\phi\to\tilde\phi$, $\mu$-a.e.~as $n\to\infty$. So, by dominated convergence theorem, it holds $\norm{\Bs_{f}^{n}\phi-\tilde\phi}_{p}\to 0$, as $n\to\infty$.

  Now, to prove the convergence in $L^{p}$ without the additional assumption that $\phi$ is bounded, let us consider a sequence of truncations of $\phi$, \ie~for each $k\in\N$ let us consider the function $\phi^{(k)}(x):=\phi(x)\ds{1}_{\phi^{-1}[-k,k]}(x)$. Observe that $\phi^{(k)}\in L^{\infty}(X,\B,\mu)$ for every $k$ and $\phi^{(k)}\to\phi$ as $k\to\infty$, $\mu$-almost everywhere and in $L^{p}$.

  So, given any $\eps>0$ there exists $k\in\N$ such that $\norm{\phi^{(k)}-\phi}_{p}<\eps/3$, and hence,  $\norm{\Bs_{f}^{n}(\phi^{(k)}-\phi)}_{p}<\eps/3$, for every $n\geq 1$, and by Corollary~\ref{cor:Birkhoff-Lp-convergence}, it holds
  \begin{displaymath}
    \norm{\tilde\phi^{(k)}-\tilde\phi}_{p}\leq \norm{\phi^{(k)}-\phi}_{p}<\frac{\eps}{3}.
  \end{displaymath}
  Then, by the previous argument, we have $\norm{\Bs_{f}^{n}\phi^{(k)}-\tilde\phi^{(k)}}_{p}<\eps/3$, whenever $n$ is sufficiently large, and so we get
  \begin{displaymath}
    \begin{split}
      \norm{\Bs_{f}^{n}\phi-\tilde\phi}_{p} &\leq \norm{\Bs_{f}^{n}\phi - \Bs_{f}^{n}\phi^{(k)}}_{p}+\norm{\Bs_{f}^{n}\phi^{(k)}-\tilde\phi^{(k)}}_{p}  + \norm{\tilde\phi^{(k)}-\tilde\phi}_{p}\\
      &\leq \frac{\eps}{3} + \frac{\eps}{3} + \frac{\eps}{3} = \eps,
    \end{split}
  \end{displaymath}
  for $n$ large enough.

  Finally, in order to prove that $\int_{x}\phi\dd\mu=\int_{X}\tilde\phi\dd\mu$, it is enough to notice that if $\phi\in L^{p}(X,\B,\mu)$, then $\phi\in L^{1}(X,\B,\mu)$ and it holds
  \begin{displaymath}
    \begin{split}
      \abs{\int_{X}\phi\dd\mu - \int_{X}\tilde\phi\dd\mu} &= \abs{\int_{X}\Bs_{f}^{n}\phi\dd\mu - \int_{X}\phi\dd\mu} \\
      &\leq \int_{X}\abs{\Bs_{f}^{n}\phi-\tilde\phi}\dd\mu=\norm{\Bs_{f}^{n}\phi-\tilde\phi}_{1}\to 0,
    \end{split}
  \end{displaymath}
  as $n\to+\infty$.
\end{proof}


\begin{exercise}
  Show that Corollary~\ref{cor:Birkhoff-Lp-convergence} is false for $p=\infty$.
\end{exercise}


\begin{corollary}[Birkhoff theorem for continuous functions]\label{cor:Birkhoff-thm-continuous-setting}
  Let $X$ be compact metric space and $f\colon X\carr$ be a continuous map. Then there exists a Borel subset $X'\subset X$ such that $\mu(X')=1$, for every $\mu\in\M(f)$ and the following limit exists for every $x\in X'$ and every $\phi\in C^{0}(X,\R)$:
  \begin{displaymath}
    \tilde\phi(x):=\lim_{n\to+\infty} \frac{1}{n}\sum_{j=0}^{n-1}\phi\big(f^{j}(x)\big).
  \end{displaymath}
\end{corollary}

\begin{proof}
  Since $X$ is compact, the Banach space $C^{0}(X,\R)$ is separable, \ie~there exists a sequence ${(\phi_{k})}_{k\geq 1}$ which is dense in $C^{0}(X,\R)$. For each $k\in\N$ we define
  \begin{displaymath}
    X_{k}:=\left\{x\in X : \lim_{n\to+\infty}\frac{1}{n}\sum_{j=0}^{n-1}\phi_{k}\big(f^{j}(x)\big) \text{ exists}\right\}.
  \end{displaymath}
  The set $X_{k}$ is clearly a Borel set and, by Birkhoff theorem (Theorem~\ref{thm:Birkhoff-ergodic-thm}) we know that $\mu(X_{k})=1$, for every $k\in\N$ and all $\mu\in\M(f)$. So, if we define
  \begin{displaymath}
    X':=\bigcap_{k\in\N} X_{k},
  \end{displaymath}
  it clearly holds $\mu(X')=1$, for every $\mu\in\M(X)$.

  Then, let us consider an arbitrary point $x\in X'$ and a positive real number $\eps>0$. Given any $\phi\in C^{0}(X,\R)$, by denseness of the sequence $(\phi_{k})$ we know there exists $k\in\N$ such that $\norm{\phi-\phi_{k}}_{0}<\eps/2$. So we have
  \begin{displaymath}
    \begin{split}
      0&\leq \limsup_{n}\frac{1}{n}\sum_{j=0}^{n-1}\phi\big(f^{j}(x)\big) - \liminf_{n}\frac{1}{n}\sum_{j=0}^{n-1}\phi\big(f^{j}(x)\big) \\
       & \leq \abs{\limsup_{n}\frac{1}{n}\sum_{j=0}^{n-1}\phi\big(f^{j}(x)\big) - \limsup_{n}\frac{1}{n}\sum_{j=0}^{n-1}\phi_{k}\big(f^{j}(x)\big)} \\
       &\qquad+ \abs{\liminf_{n}\frac{1}{n}\sum_{j=0}^{n-1}\phi_{k}\big(f^{j}(x)\big)-\liminf_{n}\frac{1}{n}\sum_{j=0}^{n-1}\phi\big(f^{j}(x)\big)} \\
      &<\frac{\eps}{2}+\frac{\eps}{2} = \eps.
    \end{split}
  \end{displaymath}
  Therefore, the  $lim_{n}\dfrac{1}{n}\sum_{j=0}^{n-1}\phi(f^{j}(x))$ exists and the result is proved.
\end{proof}


\subsection{Some important exercises from~\cite[\S3.2.4]{VianaOliveiraFoundErgTh}}
\label{sec:some-import-exerc-3.2.4}

3.2.3, 3.2.4, 3.2.5.


\section{Ergodicity}\label{sec:ergodicity}

Let $(X,\B,\mu)$ be a probability space and $f\colon X\carr$ be a measure preserving map. We say that $(f,\mu)$ is \textbf{ergodic} when given any $A\in\B$ such that $f^{-1}(A)=A$, either $\mu(A)=0$ or $\mu(A)=1$.

Ergodic systems satisfy the following properties:

\begin{theorem}\label{thm:ergodicty-equivalent-definitions}
  Let $(X,\B,\mu)$ be a probability space and $f\colon X\carr$ be a measure preserving map. Then the following statements are equivalent:
  \begin{enumerate}
    \item $(f,\mu)$ is ergodic;
    \item given any $B\in\B$ such that $\mu\big(B\triangle f^{-1}(B)\big)=0$, then $\mu(B)\in\{0,1\}$;
    \item if $\phi\in X\to\R$ is a measurable function such that
          \begin{displaymath}
            \phi\big(f(x)\big)=\phi(x), \quad\text{for }\mu\text{-a.e. }x\in X,
          \end{displaymath}
          then there exists $c\in\R$ such that $\phi(x)=c$, for $\mu$-a.e. $x\in X$;
    \item for every $A,B\in\B$ it holds
          \begin{displaymath}
            \frac{1}{n}\sum_{j=0}^{n-1}\mu\big(f^{-j}(A)\cap B) \to\mu(A)\mu(B),
          \end{displaymath}
          as $n\to+\infty$;
    \item given $p,q\in (1,\infty)$ such that $\dfrac{1}{p}+\dfrac{1}{q}=1$, and $\phi\in L^{p}(X,\B,\mu)$ and $\psi\in L^{q}(X,\B,\mu)$, it holds
          \begin{displaymath}
            \frac{1}{n}\sum_{j=0}^{n-1}\int_{X}(\phi\circ f^{j})\psi\dd\mu \to\int_{X}\phi\dd\mu \int_{X}\psi\dd\mu,
          \end{displaymath}
          as $n\to+\infty$.
  \end{enumerate}
\end{theorem}

\begin{exercise}\label{exer:essentially-invariant-set-implies-invariant}
  Let $(X,\B,\mu)$ be a probability space and $f\colon X\carr$ be measure preserving map. Let $B\in\B$ be a measurable set such that $\mu\big(B\triangle f^{-1}(B)\big)=0$. Show that there exists $A\in\B$ satisfying $f^{-1}(A)=A$ and $\mu(A\triangle B)=0$.
\end{exercise}

\begin{proof}[Proof of Theorem~\ref{thm:ergodicty-equivalent-definitions}]
  Clearly \emph{(2)} implies \emph{(1)}. On the other hand, by Exercise~\ref{exer:essentially-invariant-set-implies-invariant} we know that \emph{(1)} implies \emph{(2)}.

  If $(f,\mu)$ is not ergodic, then there exists $A\in\B$ with $f^{-1}(A)=A$ and $0<\mu(A)<1$. Then, $\ds{1}_{A}\circ f= \ds{1}_{f^{-1}(A)} = \ds{1}_{A}$ and $\ds{1}_{A}$ is not constant $\mu$-a.e. So, \emph{(3)} implies \emph{(1)}. On the other hand, if $\phi\big(f(x)\big)=\phi(x)$ for $\mu$-a.e. point $x\in X$ and $\phi$ is not constant $\mu$-almost everywhere, then there exists $r\in\R$ such that $B:=\phi^{-1}[-\infty,r]$ satisfies $0<\mu(B)<1$, with $\mu\big(B\triangle f^{-1}(B)\big)=0$. Thus, \emph{(2)} implies \emph{(3)}.

  Then, if \emph{(4)} holds and $A\in\B$ is $f$-invariant, then we have
  \begin{displaymath}
    \mu(A)=\frac{1}{n}\sum_{j=0}^{n-1}\mu\big(A\cap f^{-1}(A)\big) \to {\mu(A)}^{2},\quad\text{as }n\to+\infty.
  \end{displaymath}
  So $\mu(A)\in\{0,1\}$ and \emph{(1)} holds.

  On the other hand, assuming \emph{(5)} we can easily prove \emph{(4)} taking $\phi:=\bb{1}_{A}$ and $\psi:=\bb{1}_{B}$.

  Finally, let us prove that \emph{(1)} implies \emph{(5)}. To do that, let us star assuming $\phi\in L^{\infty}(X,\B,\mu)$ and $\psi\in L^{q}(X,\B,\mu)$. Then, by Birkhoff ergodic theorem (Theorem~\ref{thm:Birkhoff-ergodic-thm}), and since \emph{(1)} implies \emph{(3)}, we know that
  \begin{displaymath}
    \frac{1}{n}\sum_{i=0}^{n-1}\phi\big(f^{i}(x)\big)\to\int_{X}\phi\dd\mu, \quad\text{as }n\to+\infty,
  \end{displaymath}
  and for $\mu$-a.e. $x\in X$, and we have
  \begin{displaymath}
    \abs{\frac{1}{n}\sum_{i=0}^{n-1}\phi\big(f^{i}(x)\big)\psi(x)} \leq \norm{\phi}_{\infty}\abs{\psi(x)}, \quad\forall n\in\N,
  \end{displaymath}
  where $\abs{\psi}$ is an integrable function. So, by dominated converge theorem we conclude that
  \begin{displaymath}
    \lim_{n\to+\infty}\frac{1}{n}\sum_{i=0}^{n-1}\int_{x}(\phi\circ f^{i})\psi\dd\mu = \left(\int_{x}\phi\dd\mu\right)\left(\int_{X}\psi\dd\mu\right),
  \end{displaymath}
  we wanted to prove. The general case, where $\phi\in L^{p}(X,\B,\mu)$ can be proved taking the sequence of truncations of $\phi$ given by $\phi_{k}:=\bb{1}_{[\abs{\phi}\leq k]}\phi\in L^{\infty}(X,\B,\mu)$, for every $k\in\N$. Then we have $\phi_{k}\to\phi$ in $L^{p}$ and $\mu$-almost everywhere. Then, given any $\eps>0$ there is $k\in\N$ such that $\norm{\phi-\phi_{k}}_{p}<\eps$ and, by the previous argument, there is $N\in\N$ such that
  \begin{displaymath}
    \abs{\frac{1}{n}\sum_{i=0}^{n-1}\int_{X}(\phi\circ f^{i})\psi\dd\mu - \int_{X}\psi\dd\mu\int_{X}\psi\dd\mu} < \eps, \quad\forall n\geq N.
  \end{displaymath}
  Hence we have
  \begin{displaymath}
    \begin{split}
      &\abs{\frac{1}{n}\sum_{i=0}^{n-1}\int_{X}(\phi\circ f^{i})\psi\dd\mu - \int_{X}\phi\dd\mu\int_{X}\psi\dd\mu} \\
      & \leq \abs{\frac{1}{n}\sum_{i=0}^{n-1}\int_{X}\big((\phi-\phi_{k})\circ f^{i}\big)\psi\dd\mu} + \abs{\frac{1}{n}\sum_{i=0}^{n-1}\int_{X}(\phi_{k}\circ f^{i})\psi\dd\mu - \int_{X}\phi_{k}\dd\mu\int_{X}\psi\dd\mu} \\
      &\qquad + \abs{\int_{X}(\phi_{k}-\phi)\dd\mu\int_{X}\psi\dd\mu} \leq \norm{\phi-\phi_{k}}_{p}\norm{\psi}_{q} + \eps + \norm{\phi_{k}-\phi}_{1}\norm{\psi}_{1} \\
      &\leq \eps\norm{\psi}_{q} +\eps + \norm{\phi_{k}-\phi}_{p}\norm{\psi}_{q} \leq \eps \big(2\norm{\psi}_{q} + 1\big),
    \end{split}
  \end{displaymath}
  for every $n\geq N$.
\end{proof}

We shall need a little improvement of condition \emph{(3)} of Theorem~\ref{thm:ergodicty-equivalent-definitions}:

\begin{exercise}
  Let $(X,\B,\mu)$ be a probability space and $f\colon X\carr$ be a measure preserving map, and let $\mc{A}\subset 2^{X}$ be an algebra such that $\B$ is the $\sigma$-algebra generated by $\mc{A}$ and suppose that
  \begin{displaymath}
    \frac{1}{n}\sum_{i=0}^{n-1}\mu\big(f^{-i}(A)\cap B\big) \to\mu(A)\mu(B), \quad\text{as } n\to+\infty,
  \end{displaymath}
  for every $A,B\in\mc{A}$. Then, show that $(f,\mu)$ is ergodic.
\end{exercise}



\subsection{Examples of ergodic systems}\label{sec:examples-ergodic-systems}

\subsubsection{Circle rotations}\label{sec:circle-rotations-ergodicity}

Given any $\alpha\in\R$, we define the rotation $T_{\alpha}\colon \R/\Z\to\R/\Z$ by
\begin{displaymath}
  T_{\alpha}(x+\Z) := x+\alpha +\Z, \quad\forall x\in\R.
\end{displaymath}

Since $\R/\Z$ is a compact group, its Haar measure $\lambda$ is $T_{\alpha}$-invariant for every $\alpha\in\R$.

\begin{exercise}
  Let $\pi\colon\R\to\R/\Z$ be a the natural quotient map, \ie~$\pi(x):=x+\Z$, for every $x\in\R$. Show that $\lambda=\pi_{\star}\Big(\mathrm{Leb}\big|_{[0,1]}\Big)$.
\end{exercise}

\begin{exercise}
  Given $\alpha\in\Q$, show that $(T_{\alpha},\lambda)$ is not ergodic.
\end{exercise}

\begin{theorem}\label{thm:ergodicity-irrational-circle-rotations}\
  The system $(T_{\alpha},\lambda)$ is ergodic for every $\alpha\in\R\setminus\Q$.
\end{theorem}

\begin{proof}[Proof of Theorem~\ref{thm:ergodicity-irrational-circle-rotations}]
  Let $\phi\in L^2(\T,\lambda)$ be a $T_\alpha$-invariant function, \ie~$\phi(T_\alpha(x))=\phi(x)$, for $\lambda$-a.e. $x\in\T$. Considering the Fourier series of $\phi$ and $\phi\circ T_\alpha$, we get
  \begin{displaymath}
    \phi(T_\alpha(x))=\sum_{k\in\Z}\hat\phi_k e^{2\pi ik(x+\alpha)} = \sum_{k\in\Z} \hat\phi_ke^{2\pi ikx} = \phi(x),
  \end{displaymath}
  for $\lambda$-a.e. $x\in\T$.

  That means that
  \begin{displaymath}
    \hat\phi_ke^{2\pi ik\alpha}=\hat\phi_k, \quad\forall k\in\Z,
  \end{displaymath}
  and since $\alpha\in\R\setminus\Q$, we get that $k\alpha\in\Z$ if and only if $k=0$. So, $\hat\phi_k=0$ for every $k\in\Z\setminus\{0\}$ and we conclude, $\phi$ is a constant function.
\end{proof}

\subsubsection{Expanding maps on the interval}\label{sec:expanding-maps-interval-ergodicity}

Given any $k\in\N$ with $k\geq 2$, consider the map $f\colon [0,1]\carr$ given by
\begin{displaymath}
  f(x):=kx-\floor{kx},  \quad\forall x\in [0,1].
\end{displaymath}

We have shown in \S\ref{sec:expand-maps-interval} that the measure $\lambda:=\mathrm{Leb}\big|_{[0,1]}$ is $f$-invariant. Moreover, we have
\begin{theorem}\label{thm:expanding-maps-interval-ergodicity}
  The system $(f,\lambda)$ is ergodic.
\end{theorem}

\subsubsection{Bernoulli shifts}\label{sec:Bernoulli-shifts}


\subsection{Some properties of ergodic measures}\label{sec:ergodic-measures-extreme-points}

Let $(X,\B)$ be a measurable space and $f\colon X\carr$ be a measurable map. We first observe that ergodic measures are ``minimal elements'' with respect to the preorder given by absolute continuity among invariant measures:

\begin{proposition}\label{pro:ergodic-measures-minimal-elements-absolute-continuity}
  If $\mu$ and $\nu$ are two $f$-invariant probability measures such that $\mu$ is ergodic and $\nu\ll\mu$ (\ie~$\nu(A)=0$, for every $A\in\B$ such that $\mu(A)=0$), then $\nu=\mu$.
\end{proposition}

\begin{proof}
  Given an arbitrary set $A\in\B$, by Birkhoff ergodic theorem and since $\mu$ is ergodic, we get that there exists a measure $R_{A}\in\B$ with $\mu(R_{A})=1$ such that
  \begin{displaymath}
    \frac{1}{n}\sum_{i=0}^{n-1}\ds{1}_{A}\circ f^{j}(x) \to \mu(A), \quad\text{as }n\to+\infty,\ \forall x\in R_{A}.
  \end{displaymath}
  Since $\nu\ll\mu$ and $\mu(X\setminus R_{A})=0$, we get $\nu(X\setminus R_{A})=0$. So the above Birkhoff averages converge to $\mu(A)$ for $\nu$-a.e. $x\in\R$ and consequently, $\mu(A)=\nu(A)$. Since $A$ was arbitrary, we get $\mu=\nu$.
\end{proof}

We get the following
\begin{corollary}\label{cor:ergodic-measures-as-extreme-points}
  Let $(X,\B)$ be a measurable space and $f\colon X\carr$ be a measurable map. If $\mu$ denotes an $f$-invariant probability measure, then $\mu$ is not ergodic if and only there exist two $f$-invariant measures $\nu_{1}$ and $\nu_{2}$ with $\nu_{1}\neq\nu_{2}$ and $t\in (0,1)$ such that
  \begin{displaymath}
    \mu = t\nu_{1} + (1-t)\nu_{2}.
  \end{displaymath}
\end{corollary}

\begin{proof}
  Let us start assuming $\mu$ is not ergodic. Then there is $A\in\B$ with $0<\mu(A)<1$. So, let us define $\nu_{!}(E):=\dfrac{1}{\mu(A)}\mu(A\cap E)$ and $\nu_{2}:=\dfrac{1}{\mu(X\setminus A)}\mu(E\setminus A)$, for every $E\in\B$, and it clearly holds
  \begin{displaymath}
    \mu = \mu(A)\nu_{1} + \mu(X\setminus A)\nu_{2} = \mu(A)\nu_{1}+\big(1-\mu(A)\big)\nu_{2}.
  \end{displaymath}

  Reciprocally, let us suppose that $\mu$ is ergodic and there are $f$-invariant measures $\nu_{1}$ and $\nu_{2}$ and $t\in(0,1)$ such that $\mu=t\nu_{1}+(1-t)\nu_{2}$. Then it clearly holds $\nu_{1}\ll\mu$ and $\nu_{2}\ll\mu$. By Proposition~\ref{pro:ergodic-measures-minimal-elements-absolute-continuity}, we know that $\mu=\nu_{1}=\nu_{2}$.
\end{proof}


\subsection{Some important exercises from~\cite[\S4.1.3]{VianaOliveiraFoundErgTh}}
\label{sec:some-import-exerc-4.1.3}

4.1.1, 4.1.2, 4.1.3, 4.1.4, 4.1.5, 4.1.7.

\subsection{Some important exercises from~\cite[\S4.2.7]{VianaOliveiraFoundErgTh}}
\label{sec:some-import-exerc-4.2.7}

4.2.4, 4.2.5, 4.2.6 (be careful!), 4.2.8, 4.2.11, 4.2.12.



\section{Unique ergodicity}\label{sec:unique-ergodicity}


\section{Ergodic decomposition}\label{sec:ergod-decomposition}


\subsection{Partitions of measure spaces}\label{sec:partitions-measure-spaces}

Let $(X,\B,\mu)$ be an arbitrary probability space. A \textbf{partition} of $X$ is a family $\xi\subset\B$ such that $C\cap D=\emptyset$, for every $C,D\in\xi$ with $C\neq D$ and
\begin{displaymath}
  \mu\bigg(\bigsqcup_{C\in\xi} C\bigg) = 1.
\end{displaymath}
The elements of $\xi$ are called \textbf{atoms} of $\xi$. Given a point $x\in\bigsqcup_{C\in\xi} C$, we define $\xi(x)\in\xi$ as the only atom of $\xi$ that contains the point $x$. So, we can define the natural \emph{projection map} $\pi^{\xi}\colon X\to\xi$ by $\pi^{\xi}(x):=\xi(x)$ and notice that $\pi^{\xi}$ is well defined $\mu$-almost everywhere.

We can use the natural projection map $\pi^{\xi}$ to endow the set $\xi$ with a measure space structure. Indeed, we shall always assume the partition $\xi$ as the supporting set of the measure space $(\xi,\B_{\xi},\hat\mu^{\xi})$ where $\B_{\xi}:=\pi^{\xi}_{\star}\B$ and $\hat\mu^{\xi}:=\pi^{\xi}_{\star}\mu$ denote the push-forward $\sigma$-algebra of $\B$ and the measure $\mu$ by $\pi^{\xi}$, respectively.

We can define a partial order relation among partitions of a measure space. Given two partitions $\xi,\eta$ of the measure space $(X,\B,\mu)$, we say that $\xi$ is \textbf{finer} than $\eta$, or that $\eta$ is \textbf{coarser} that $\xi$, and we write $\eta\leq\xi$, when there is a measurable subset $X'\subset X$ with $\mu(X')=1$ and such that
\begin{displaymath}
  \xi(x)\cap X'\subset \eta(x)\cap X', \quad\forall x\in \bigg(\bigsqcup_{C\in\xi} C\bigg)\cap \bigg(\bigsqcup_{D\in\eta} D\bigg).
\end{displaymath}
We say that $\xi$ and $\eta$ coincide or are \textbf{equal} when $\eta\leq \xi$ and $\xi\leq\eta$.

If we write $\mc{P}_{X}$ to denote the set of all partitions of the measure space $(X,\B,\mu)$, we know that $(\mc{P}_{X},\leq)$ is a partially order set.

\begin{exercise}\label{exer:inf-sup-well-defined-on-partitions}
  Show that the notions of \emph{infimum} and \emph{supremum} are well defined for countable subsets on the partially order set $(\mc{P}_{X},\leq)$. More precisely, show that given any countable family of partitions $\Xi\subset\mc{P}_{X}$, there exists two partitions $\inf\Xi,\sup\Xi\in\mc{P}_{X}$ such that
  \begin{displaymath}
    \eta\leq\inf\Xi\leq \xi\leq \sup\Xi\leq\eta', \quad\forall \xi\in\Xi,
  \end{displaymath}
  and every $\eta,\eta'\in\mc{P}_{X}$ such that $\eta\leq \xi'\leq \eta'$, for every $\xi'\in\mc{P}_{X}$.
\end{exercise}

We can use the notions of infimum and supremum on $\mc{P}_{X}$ to define the operators $\wedge$ and $\vee$: given any countable (or finite) family $\Xi\subset\mc{P}_{X}$, we define
\begin{displaymath}
  \bigwedge_{\xi\in\Xi} \xi :=\sup\left\{\eta\in\mc{P}_{X} : \eta\leq\xi,\ \forall \xi\in\Xi\right\},
\end{displaymath}
and
\begin{equation}\label{eq:partitions-upper-limit-def}
  \bigvee_{\xi\in\Xi} \xi :=\inf\left\{\eta\in\mc{P}_{X} : \xi\leq\eta,\ \forall\xi\in\Xi\right\}.
\end{equation}


\subsection{Disintegration of measures}\label{sec:disintegrations-measures}

Given a probability space $(X,\B,\mu)$ and a partition $\xi\in\mc{P}_{X}$, a \textbf{disintegration} of $\mu$ with respect to $\xi$ is family $\{\mu_{C} : C\in\xi\}$ satisfying the following properties:
\begin{enumerate}[(i)]
  \item $\mu_{C}$ is a probability measure on the measurable space $(X,\B)$, for every $C\in\xi$;
  \item $\mu_{C}(C)=1$, for $\hat\mu^{\xi}$-almost every atom $C\in\xi$;
  \item for every $E\in\B$, the map $\xi\ni C\mapsto \mu_{C}(E)\in[0,1]$ is $(\B^{\xi},\B_{\R})$-measurable;
  \item it holds
        \begin{displaymath}
          \mu(E)=\int_{\xi} \mu_{C}(E)\dd\hat\mu^{\xi}(C), \quad\forall E\in\B.
        \end{displaymath}
\end{enumerate}

Let us observe that one can pull back the partition and push forward the disintegration:

\begin{exercise}\label{exer:push-forward-disint-pull-back-partitions}
  Let $(X,\B,\mu)$ be a probability space, $\xi\in\mc{P}_{X}$ be a partion and $f\colon X\carr$ be a measure preserving map. Then, show that $f^{\star}\xi:=\{f^{-1}(C) : C\in\xi\}$ is a partition of $(X,\B,\mu)$. On the other hand, if $\{\mu_{D} : D\in f^{\star}\xi\}$ is a disintegration of $\mu$ with respect to the partition $f^{\star}\xi$, show that $\{\mu_{C}:=f_{\star}\mu_{f^{-1}(C)} :  C\in\xi\}$ is a disintegration of $\mu$ with respect to $\xi$.
\end{exercise}



First we shall prove that every partition admits at most one disintegration:

\begin{theorem}[Uniqueness of disintegrations]\label{thm:unique-disintegration}
  If $(X,\B,\mu)$ is a Borel space (see \S\ref{sec:measure-theory}), $\xi\in\mc{P}_{X}$ is a partition, and $\{\mu_{C} : C\in\xi\}$ and $\{\nu_{C} : C\in\xi\}$ are two disintegrations of $\mu$ with respect to $\xi$, then it holds
  \begin{displaymath}
    \mu_{C}=\nu_{C}, \quad\text{for } \hat\mu^{\xi}\text{-a.e. } C\in\xi.
  \end{displaymath}
\end{theorem}

\begin{proof}
  Since $(X,\B,\mu)$ is a Borel space, there exists a countable algebra $\ms{A}\subset\B$ that generates $\B$.
  For each $A\in\ms{A}$, let us consider the sets
  \begin{displaymath}
    \ms{M}_{A}:=\{C\in\xi : \mu_{C}(A)<\nu_{C}(A)\},\quad\text{and } \ms{N}_{A}:=\{C\in\xi : \mu_{C}(A)>\nu_{C}(A)\}.
  \end{displaymath}

  Then, since $\{\mu_{C} : C\in\xi\}$ is a disintegration of $\mu$ with respect to $\xi$, we have
  \begin{displaymath}
    \begin{split}
      \mu\big(A\cap{\pi^{\xi}}^{-1}(\ms{M}_{A})\big) &= \int_{\xi} \mu_{C}\big(A\cap{\pi^{\xi}}^{-1}(\ms{M}_{A})\big)\dd\hat\mu^{\xi}(C)\\
      &=\int_{\ms{M}_{A}} \mu_{C}(A)\dd\hat\mu^{\xi}(C).
    \end{split}
  \end{displaymath}

  Analogously, one can show that
  \begin{displaymath}
    \mu\big(A\cap{\pi^{\xi}}^{-1}(\ms{M}_{A})\big)=\int_{\ms{M}_{A}} \nu_{C}(A)\dd\hat\mu^{\xi}(C).
  \end{displaymath}

  Since $\mu_{C}(A)<\nu_{C}(A)$, for all $C\in\ms{M}_{A}$, we get $\hat\mu^{\xi}(\ms{M}_{A})=0$. Analogously, we can show that $\hat\mu^{\xi}(\ms{N}_{A})=0$ and hence, $\mu_{C}=\nu_{C}$, for $\hat\mu^{\xi}$-almost every $C\in\xi$.
\end{proof}

Theorem~\ref{thm:unique-disintegration} can be used to show that certain partitions do not admit any disintegration of the measure:

\begin{example}\label{exam:disintegration-non-existence}
  Let $\alpha$ be an arbitrary real irrational number and $f\colon\R/\Z\carr$ be the corresponding rotation $f(x+\Z):=x+\alpha+\Z$, for every $x\in\R$. Let $\lambda$ denote the Haar measure on the circle $\R/\Z$ and consider the partition $\xi$ of the probability space $(\R/\Z,\B_{\R/\Z},\lambda)$ given by the orbits of $f$, \ie~we have
  \begin{displaymath}
    \xi:=\left\{\{f^{n}(x)\in\R/\Z : n\in\Z\} : x\in\R/\Z\right\}.
  \end{displaymath}

  Then there is no disintegration of $\lambda$ with respect to the partion $\xi$.
\end{example}

Now we can state the main result of this section:

\begin{theorem}[Ergodic decomposition theorem]\label{thm:ergodic-decomposition}
  Let $(X,\B,\mu)$ be a Borel space and $f\colon X\carr$ be a measure preserving map. Then there exists a unique partition $\xi_{f}\in\mc{P}_{X}$, called the \textbf{ergodic partition of $f$}, such $\xi_{f}$ admits a disintegration $\{\mu_{C} : C\in\xi_{f}\}$ and the following conditions hold:
  \begin{enumerate}
    \item $f^{\star}\xi_{f}=\xi_{f}$, \ie~$f^{-1}(C)=C$, for $\hat\mu^{\xi_{f}}$-almost every atom $C\in\xi_{f}$;
    \item $\mu_{C}$ is an ergodic $f$-invariant measure, for $\hat\mu^{\xi_{f}}$-almost every atom $C\in\xi_{f}$.
  \end{enumerate}
\end{theorem}

The proof of Theorem~\ref{thm:ergodic-decomposition} will be postpone and get as a consequence of Rokhlin decomposition theorem.

\subsection{Measurable partitions}\label{sec:measurable-partitions}

By Example~\ref{exam:disintegration-non-existence} we know that some partitions do not admit disintegrations. In this section we shall characterize the family of partitions admitting disintegrations.

Given a probability space $(X,\B,\mu)$ we say that a partition $\xi\in\mc{P}_{X}$ is a \textbf{measurable partition} when there exists a sequence ${(\xi_{n})}_{n\geq 1}\subset\mc{P}_{X}$ of finite partitions such that
\begin{displaymath}
  \xi=\bigvee_{n\geq 1}\xi_{n},
\end{displaymath}
where the equality among partitions is given in \S\ref{sec:partitions-measure-spaces} and the operator $\bigvee$ is given by~\eqref{eq:partitions-upper-limit-def}

\begin{exercise}
  Let $(X,\B,\mu)$ be a probability space, $M$ be separable metric space and $f\colon X\to M$ be a measurable map. Then the partition
  \begin{displaymath}
    \xi_{f}:=\left\{f^{-1}(y) : y \in M\right\}\in\mc{P}_{X}
  \end{displaymath}
  is a measurable partition.
\end{exercise}

Then we have the main result about measurable partitions:

\begin{theorem}[Rokhlin theorem]\label{thm:rokhlin-disintegration}
  Let $(X,\B,\mu)$ be a Borel space and $\xi\in\mc{P}_{X}$ be measurable partition. Then there exits a disintegration of $\mu$ with respect to $\xi$.
\end{theorem}

We postpone the proof of Theorem~\ref{thm:rokhlin-disintegration}. Now we shall prove the Ergodic decomposition theorem assuming Theorem~\ref{thm:rokhlin-disintegration}.

\begin{proof}[Proof of Theorem~\ref{thm:ergodic-decomposition}]
  Since $(X,\B,\mu)$ is a Borel space, there exists a countable algebra $\ms{A}\subset\B$ that generates the $\sigma$-algebra $\B$. For each $A\in\ms{A}$, let us consider the function $\tau_{A}\colon X\to [0,1]$ given by
  \begin{displaymath}
    \tau_{A}(x):=\lim_{n\to+\infty}\frac{1}{n}\sum_{i=0}^{n-1}\ds{1}_{A}\big(f^{i}(x)\big),
  \end{displaymath}
  whenever the limit exists. By Birkhoff ergodic theorem we know that $\tau_{A}$ is well defined $\mu$-almost everywhere.

  Since $\ms{A}$ is countable, we get the set $X_{0}:=\{x\in X : \tau_{A}(x) \text{ is well defined, }\forall A\in\ms{A}\}$ is a measurable set with $\mu(X_{0})=1$.

  Then we define the partition $\xi_{f}\in\mc{P}_{X}$ as the set of equivalence classes on $X_{0}$ such that $x\sim y$ if and only if $\tau_{A}(x)=\tau_{A}(y)$, for every $A\in\ms{A}$.

  \begin{exercise}
    Show that $\xi_{f}$ is a measurable partition.
  \end{exercise}

  Then, by Rokhlin theorem (Theorem~\ref{thm:rokhlin-disintegration}) we know that there is a disintegration of $\mu$ with respect to $\xi_{f}$ $\{\mu_{C} : C\in\xi_{f}\}$.

  By the very same definition of functions $\tau_{A}$ we know that $f^{-1}(X_{0})=X_{0}$ and  $\tau_{A}(x)=\tau_{A}\circ f(x)$, for every $x\in X_{0}$. So, it clearly holds $f^{-1}(C)=C$, for every $C\in\xi_{f}$. Thus, by uniqueness of disintegration (Theorem~\ref{thm:unique-disintegration}), we conclude that $f_{\star}\mu_{C}=\mu_{C}$, for $\hat\mu^{\xi_{f}}$-almost every atom $C\in\xi_{f}$.

  So, it just remains to prove that the measure $\mu_{C}$ is ergodic for $f$, for every atom $C\in\xi_{f}$. To do that, let us fix an atom $C\in\xi_{f}$ and observe that to prove the ergodicity of $\mu_{C}$ is enough to show that $\tau_{E}(x)=\mu_{C}(E)$, for every $E\in\B$ and $\mu_{C}$-a.e. $x\in C$. Then, let us define the family
  \begin{displaymath}
    \B_{C}:=\left\{E\in\B : \tau_{E}(x)=\mu_{C}(E),\ \text{for }\mu_{C}\text{-a.e. }x\in C\right\}.
  \end{displaymath}
  By our definition it clearly holds $\ms{A}\subset\B_{C}$. On the other hand, if $E_{1},E_{2}\in\B_{C}$, and $E_{1}\subset E_{2}$, then $E_{2}\setminus E_{1}\in\B_{C}$ because
  \begin{displaymath}
    \tau_{E_{2}\setminus E_{1}}(x)=\tau_{E_{2}}(x)-\tau_{E_{1}}(x) = \mu_{C}(E_{2}\setminus E_{1}),
  \end{displaymath}
  for $\mu_{C}$-almost every $x\in C$; and if $E_{1},E_{2},\ldots\in\B_{C}$ is a family of pairwise disjoint sets, the it clearly holds
  \begin{displaymath}
    \tau_{\bigsqcup_{j}E_{j}}(x) = \sum_{j}\tau_{E_{j}}(x)= \mu_{C}\bigg(\bigsqcup_{j} E_{j}\bigg).
  \end{displaymath}

  Then, the ergodicity of $\mu_{C}$ is consequence of the following
  \begin{exercise}
    A family $\mc{C}\subset 2^{X}$ is called a \textbf{monotone class} when given any sequence $A_{1}\subset A_{2}\subset\ldots$ and $B_{1}\supset B_{2}\supset \ldots$, with $A_{n},B_{n}\in\mc{C}$, for every $n\in\N$, it holds $\bigcup_{n} A_{n}, \bigcap_{n} B_{n}\in \mc{C}$.

    Then, if $\mc{C}$ is a monotone class and contains an algebra $\ms{A}$, then $\mc{C}$ contains the $\sigma$-algebra $\sigma(\ms{A})$ generated by $\ms{A}$.
  \end{exercise}
\end{proof}


\subsection{Some important exercises from~\cite[\S5.1.5 and \S5.2.4]{VianaOliveiraFoundErgTh}}
\label{sec:some-import-exerc-5.1.5}

5.1.1, 5.1.2, 5.1.3, 5.1.6, 5.1.7.

5.2.1.




% --------- References ----------------------------------------------------------

\bibliographystyle{amsalpha}
\bibliography{references.bib}

% ------------------------------------------------------------------------------

\end{document}
